% Copyright (C) 2012 Raniere Silva
% 
% This file is part of 'CNPq 126874/2012-3'.
% 
% 'CNPq 126874/2012-3' is licensed under the Creative Commons
% Attribution 3.0 Unported License. To view a copy of this license,
% visit http://creativecommons.org/licenses/by/3.0/.
% 
% 'CNPq 126874/2012-3' is distributed in the hope that it will be
% useful, but WITHOUT ANY WARRANTY; without even the implied warranty of
% MERCHANTABILITY or FITNESS FOR A PARTICULAR PURPOSE.

\documentclass[12pt,a4paper]{article}
% Filename: packages.tex
% This code is part of 'CNPq 126874/2012-3'.
% 
% Description: Pacotes utilizados.
% 
% Created: 20.08.12 07:27:58 PM
% Last Change: 20.08.12 07:27:58 PM
% 
% Author: Raniere Silva, <r.gaia.cs@gmail.com>
% 
% Copyright (c) 2012, Raniere Silva. All rights reserved.
% 
% This file is license under the terms of the GNU General Public License as published by the Free Software Foundation, either version 3 of the License, or (at your option) any later version. More details at <http://www.gnu.org/licenses/>
%
%Pacotes utilizados
\usepackage[top=3cm,left=2cm,right=2cm,bottom=3cm]{geometry} %Borda das p\'{a}gina
\usepackage[utf8]{inputenc} %Ser\'{a} utilizado a codifica\c{c}ao UTF8
\usepackage[T1]{fontenc}
%\usepackage[latin1]{inputenc} % Ser\'{a} utilizado a codifica\c{c}ao latin1
\usepackage[brazil]{babel} %Ser\'{a} utilizado o idioma portugu\^{e}s
\usepackage{indentfirst} %Identa\c{c}\~{a}o de linha
\usepackage{setspace} %Espa\c{c}amento entre linas
\usepackage{amsmath, amsfonts, amssymb, amsthm} % Pacote matem\'{a}tico
\usepackage{graphicx} %Pacote para incluir figuras
\usepackage{color} %Pacote para cores
\usepackage{tikz} %Pacote para desenho de figuras
\usepackage{pgfplots} %Pacote para desenho de gr\'{a}ficos
\usepackage{subfigure} % Pacote para subfiguras
\usepackage{url} % Pacote para url
\usepackage{hyperref} % Pacote para hyperlink
\usepackage{algorithmic} % Pacote para algoritmos
\usepackage{algorithm} % Pacote para algoritmos
\usepackage{listings} % Pacote para c\'{o}digos

%Personalização
%\pgfplotsset{
%height=0.35\textheight,
%width=0.8\textwidth,
%legend style={
%at={(0.3,0.97)},
%anchor=north,
%font=\footnotesize
%},
%xlabel={},
%ylabel={},
%xtick={},
%grid=major,
%}

\hypersetup{
%colorlinks = true,
}

\floatname{algorithm}{Algoritmo}

\algsetup{linenosize=\small}
\renewcommand{\algorithmicrequire}{\textbf{Entrada:}}
\renewcommand{\algorithmicensure}{\textbf{Saída:}}
\renewcommand{\algorithmicend}{\textbf{fim}}
\renewcommand{\algorithmicif}{\textbf{se}}
\renewcommand{\algorithmicthen}{\textbf{ent\~{a}o}}
\renewcommand{\algorithmicelse}{\textbf{caso contr\'{a}rio}}
\renewcommand{\algorithmicendif}{\algorithmicend}
\renewcommand{\algorithmicfor}{\textbf{para}}
\renewcommand{\algorithmicforall}{\textbf{para todo}}
\renewcommand{\algorithmicdo}{\textbf{fa\c{c}a}}
\renewcommand{\algorithmicendfor}{\algorithmicend}
\renewcommand{\algorithmicwhile}{\textbf{enquanto}}
\renewcommand{\algorithmicendwhile}{\algorithmicend}
\renewcommand{\algorithmicrepeat}{\textbf{repita}}
\renewcommand{\algorithmicuntil}{\textbf{at\'{e}}}
\renewcommand{\algorithmicreturn}{\textbf{retorne}}
\renewcommand{\algorithmiccomment}[1]{\hspace{2em}/* #1 */}

\renewcommand{\lstlistingname}{C\'{o}digo}
\lstset{ %
% language=Octave,                % the language of the code
basicstyle=\ttfamily\small,       % the size of the fonts that are used for the code
numbers=left,                   % where to put the line-numbers
numberstyle=\footnotesize,      % the size of the fonts that are used for the line-numbers
stepnumber=5,                   % the step between two line-numbers. If it's 1, each line 
% will be numbered
numbersep=5pt,                  % how far the line-numbers are from the code
% backgroundcolor=\color{white},  % choose the background color. You must add \usepackage{color}
showspaces=false,               % show spaces adding particular underscores
showstringspaces=false,         % underline spaces within strings
showtabs=false,                 % show tabs within strings adding particular underscores
% frame=single,                   % adds a frame around the code
tabsize=4,                      % sets default tabsize to 2 spaces
captionpos=t,                   % sets the caption-position to bottom
breaklines=true,                % sets automatic line breaking
breakatwhitespace=false,        % sets if automatic breaks should only happen at whitespace
caption={\texttt{\lstname}},                 % show the filename of files included with \lstinputlisting;
% also try caption instead of title
% escapeinside={\%*}{*)},         % if you want to add a comment within your code
% morekeywords={#}            % if you want to add more keywords to the set
}

\newtheorem{defi}{Definição}
\newtheorem{prop}{Proposição}

\begin{document}
\title{Implementa\c{c}\~{a}o eficiente da heur\'{i}stica de reordenamento de Cuthill-McKee reversa}
\author{Raniere Silva \\ Orientando \and Aurelio Oliveira \\ Orientador}
\maketitle

% ver orientações em http://www.prp.unicamp.br/pibic/relatorio_final.php
% Relatório deverá conter no máximo 20 páginas, incluindo bibliografia e anexos.
\tableofcontents

% Identificação
%     Projeto
%     Bolsista / RA
%     Orientador
%     Local de execução
%     Vigência

% Introdução
%     Introdução ao assunto: deve ser bastante geral
%     Informações da literatura: tornam a introdução mais específica ao assunto
%     Colocação da questão estudada: especificar os objetivos do trabalho
%     Atividades desenvolvidas: dar uma idéia geral de como foi desenvolvido o
%     trabalho
% Copyright (C) 2012 Raniere Silva
% 
% This file is part of 'CNPq 126874/2012-3'.
% 
% 'CNPq 126874/2012-3' is licensed under the Creative Commons
% Attribution 3.0 Unported License. To view a copy of this license,
% visit http://creativecommons.org/licenses/by/3.0/.
% 
% 'CNPq 126874/2012-3' is distributed in the hope that it will be
% useful, but WITHOUT ANY WARRANTY; without even the implied warranty of
% MERCHANTABILITY or FITNESS FOR A PARTICULAR PURPOSE.

\section{Introdução}
Neste trabalho revisitou-se a heurística de reordenamento Cuthill-McKee Reversa,
proposta em 1969, que busca um reordenamento para matrizes esparsas que reduza a
largura de banda destas. O autor deste trabalho implementou a heurística
revisitada, testou com várias bibliotecas (dentre elas a ``Netlib LP'' e
``Meszaros'') e concluiu, pelo menos para os problemas testados, que a
heurística Cuthill-McKee Reversa é inferior a heurística de mínimo grau
múltiplo por gerar mais elementos não nulos na decomposição de Cholesky.

\section{Métodos de Pontos Interiores Primal-Dual}
Consideremos o Problema de Programação Linear na Forma Padrão (PPLFP):
\begin{align*}
    \text{minimizar } & c^T x \\
    \text{sujeito a } & A x = b, \\
    & x \geq 0,
\end{align*}
onde $A \in \mathbb{R}^{m \times n}$ é uma matriz de posto completo $m$ e $c$,
$b$ e $x$ são vetores colunas de dimensão apropriada. Associado a este problema
temos o problema dual (PPLFD):
\begin{align*}
    \text{maximizar } & b^T y \\
    \text{sujeito a } & A^T y + z = c, \\
    & z \geq 0,
\end{align*}
onde $y$ é um vetor coluna de dimensão $m$ de variáveis livres e $z$ é o vetor
coluna de dimensão $n$ de variáveis de folga duais. O \textit{gap} dual é dado
por $\gamma = c^T x - b^T y$ que se reduz a $\gamma = x^T z$ para pontos primais
e duais factíveis.

Os Métodos de Pontos Interiores Primais-Duais para resolver PPLFP consistem em, a
partir de uma tripla inicial $(x^0, y^0, z^0)$, construir uma
sequência de triplas, $(x^i, y^i, z^i)$, dada por
\begin{align*}
    x^{i + 1} &= x^i + \alpha^i \Delta x^i, \\
    y^{i + 1} &= y^i + \alpha^i \Delta y^i, \\
    z^{i + 1} &= z^i + \alpha^i \Delta z^i,
\end{align*}
onde $i \geq 0$ e $\alpha^i \in (0, 1]$, que tende para a tripla $(x^*, y^*, z^*)$ que é
solução de PPLFP e PPLFD. A constante $\alpha^i$ deve ser escolhida tal que $x^i,
z^i > 0$ e, frequentemente, $\alpha^i << 1$ \cite{Wright:1987:PrimalDualIP}.

A direção afim nos Métodos de Pontos Interiores Primais-Duais, $(\Delta x^i,
\Delta y^i, \Delta z^i)$,  é dada por:
\begin{align}
    \begin{bmatrix}
         0 & A^T & I \\
         A & 0 & 0 \\
         Z & 0 & X
     \end{bmatrix} \begin{bmatrix}
         \Delta x^i \\
         \Delta y^i \\
         \Delta z^i
     \end{bmatrix} &= \begin{bmatrix}
         r_d \\
         r_p \\
         r_a
     \end{bmatrix},
     \label{eq:primal_dual_intp:lin_system}
\end{align}
onde $X = \diag(x^i)$, $Z = \diag(z^i)$, $r_p = b - A x^i$, $r_d = c - A^T y^i
- z^i$, $r_a = - X Z e$ e $e$ representa o vetor de uns.

Eliminando as variáveis $\Delta z^i$ de \eqref{eq:primal_dual_intp:lin_system}
obtemos o sistema aumentado:
\begin{align}
    \begin{bmatrix}
        -D & A^T \\
        A & 0
    \end{bmatrix} \begin{bmatrix}
        \Delta x^i \\
        \Delta y^i
    \end{bmatrix} &= \begin{bmatrix}
        r_1 \\
        r_2
    \end{bmatrix},
    \label{eq:primal_dual_intp:aug_lin_system}
\end{align}
onde $D = X^{-1} Z$, $r_1 = r_d - X^{-1} r_a$ e $r_2 = r_p$, e
\begin{align*}
    \Delta z^i &= X^{-1} \left( r_a - Z \Delta x^i \right).
\end{align*}

A forma mais utilizada para resolver \eqref{eq:primal_dual_intp:aug_lin_system}
consiste em reduzir o sistema através de eliminação das variáveis $\Delta x^i$ de
forma a obter o sistema dado por
\begin{align}
    A D^{-1} A^T \Delta y^i &= A D^{-1} r_1 - r_2
    \label{eq:primal_dual_intp:lin_norm}
\end{align}
e que
\begin{align*}
    \Delta x^i &= D^{-1} \left( r_1 - A^T \Delta y^i \right).
\end{align*}

Para resolver \eqref{eq:primal_dual_intp:lin_norm} costuma-se utilizar a
fatoração de Cholesky pois $A D^{-1} A^T$ é uma matriz simétrica definida
positiva. Durante a construção da sequência de triplas $(x^i, y^i, z^i)$ apenas
$D$ se altera na matriz e por esse motivo a estrutura esparsa de $A D^{-1} A^T$ é mantida
durante todas as iterações do Método de Pontos Interiores Primal-Dual.

\section{Matrizes Esparsas}
Como visto na seção anterior, um dos passos dos Métodos de Pontos Interiores
Primal-Dual consiste em resolver um sistema linear simétrico definido positivo,
cuja estrutura esparsa é mantida inalterada durante todo o método, utilizando a fatoração de
Cholesky. Na próxima subseção introduziremos as definições de banda e envelope
de matrizes para posteriormente apresentarmos a fatoração de Cholesky.

\subsection{Banda e Envelope de Matrizes}
Mudando a notação, considere $A \in \mathbb{R}^{n \times n}$ uma matriz
simétrica definida positiva genérica com entradas $A_{ij}$. Para cada linha $i$
de $A$, $i = 1, \ldots, n$, seja
\begin{align*}
    f_i(A) &= \min\left\{ j \mid A_{ij} \neq 0 \right\},
\end{align*}
isso é, $f_i(A)$ corresponde a primeira coluna cujo elemento na linha $i$ é
diferente de zero.

Seguindo Cuthill e McKee \cite{Cuthill:1969:ReducingBandwidth}, definimos
\begin{align*}
    \beta_i(A) &= i - f_i(A), 1 \leq i \leq n, \\
    \beta(A) &= \max\left\{ \beta_i(A) \mid 1 \leq i \leq n \right\}
\end{align*}
em que $\beta_i(A)$ é a largura de banda da $i$-ésima linha de $A$ e
$\beta(A)$ é a largura de banda da matriz $A$. A banda da matriz $A$ é
definida como
\begin{align*}
    \text{Band}(A) &= \left\{ \left\{ i, j \right\} \mid 0 < i - j \leq
    \beta(A) \right\}
\end{align*}
e ilustrada na Figura~\ref{fig:exem_bandwidth}.
\begin{figure}[!htb]
    \centering
    \begin{tikzpicture}
        \matrix (A) [matrix of math nodes,%
        left delimiter  = \lbrack,%
        right delimiter = \rbrack] at (0,0)
        {%
            1 & 1 & 1 & 1 & 1 & 0 & 0 \\
            1 & 2 & 1 & 1 & 1 & 1 & 1 \\
            1 & 1 & 2 & 1 & 1 & 1 & 1 \\
            1 & 1 & 1 & 2 & 1 & 0 & 0 \\
            1 & 1 & 1 & 1 & 2 & 0 & 0 \\
            0 & 1 & 1 & 0 & 0 & 3 & 2 \\
            0 & 1 & 1 & 0 & 0 & 2 & 3 \\
        };
        \node[above, shift={(0,.5)}] at (A-1-1) {$1$};
        \node[above, shift={(0,.5)}] at (A-1-2) {$2$};
        \node[above, shift={(0,.5)}] at (A-1-3) {$3$};
        \node[above, shift={(0,.5)}] at (A-1-4) {$4$};
        \node[above, shift={(0,.5)}] at (A-1-5) {$5$};
        \node[above, shift={(0,.5)}] at (A-1-6) {$6$};
        \node[above, shift={(0,.5)}] at (A-1-7) {$7$};
        \node[right, shift={(-1.5,0)}] at (A-1-1) {$1$};
        \node[right, shift={(-1.5,0)}] at (A-2-1) {$2$};
        \node[right, shift={(-1.5,0)}] at (A-3-1) {$3$};
        \node[right, shift={(-1.5,0)}] at (A-4-1) {$4$};
        \node[right, shift={(-1.5,0)}] at (A-5-1) {$5$};
        \node[right, shift={(-1.5,0)}] at (A-6-1) {$6$};
        \node[right, shift={(-1.5,0)}] at (A-7-1) {$7$};
        % Largura de banda
        % \draw (A-1-6.north east) -- (A-2-7.north east);
        \draw (A-6-1.south west) -- (A-7-2.south west);
        % Banda
        \fill[gray, fill opacity=0.2, draw=black, dotted] (A-2-1.north west) -- (A-6-1.south west)
        -- (A-6-2.south west) -- (A-7-2.south west) -- (A-7-7.south west)
        \foreach \x in {6,5,4,3,2}{
        -- (A-\x-\x.south east) -- (A-\x-\x.south west) -- (A-\x-\x.north west)
        } -- (A-2-1.north west);

        \matrix (B) [matrix of math nodes,%
        left delimiter  = \lbrack,%
        right delimiter = \rbrack] at (8,0)
        {%
           10 &  0 &  2 &  3 &  0 &  0 &  0 \\
            0 & 10 &  2 &  3 &  0 &  0 &  0 \\
            2 &  2 & 10 &  0 &  4 &  0 &  0 \\
            3 &  3 &  0 & 10 &  4 &  0 &  0 \\
            0 &  0 &  4 &  4 & 10 &  2 &  2 \\
            0 &  0 &  0 &  0 &  2 & 10 &  0 \\
            0 &  0 &  0 &  0 &  2 &  0 & 10 \\
        };
        \node[above, shift={(0,.5)}] at (B-1-1) {$1$};
        \node[above, shift={(0,.5)}] at (B-1-2) {$2$};
        \node[above, shift={(0,.5)}] at (B-1-3) {$3$};
        \node[above, shift={(0,.5)}] at (B-1-4) {$4$};
        \node[above, shift={(0,.5)}] at (B-1-5) {$5$};
        \node[above, shift={(0,.5)}] at (B-1-6) {$6$};
        \node[above, shift={(0,.5)}] at (B-1-7) {$7$};
        \node[right, shift={(-1.5,0)}] at (B-1-1) {$1$};
        \node[right, shift={(-1.5,0)}] at (B-2-1) {$2$};
        \node[right, shift={(-1.5,0)}] at (B-3-1) {$3$};
        \node[right, shift={(-1.5,0)}] at (B-4-1) {$4$};
        \node[right, shift={(-1.5,0)}] at (B-5-1) {$5$};
        \node[right, shift={(-1.5,0)}] at (B-6-1) {$6$};
        \node[right, shift={(-1.5,0)}] at (B-7-1) {$7$};
        % Largura de banda
        % \draw (B-1-4.north east) -- (B-4-7.north east);
        \draw (B-4-1.south west) -- (B-7-4.south west);
        % Banda
        \fill[gray, fill opacity=0.2, draw=black, dotted] (B-1-1.south west) -- (B-4-1.south west)
        -- (B-4-2.south west) -- (B-5-2.south west) -- (B-5-3.south west)
        -- (B-6-3.south west) -- (B-6-4.south west) -- (B-7-4.south west)
        -- (B-7-7.south west)
        \foreach \x in {6,5,4,3,2}{
        -- (B-\x-\x.south east) -- (B-\x-\x.south west) -- (B-\x-\x.north west)
        } -- (B-1-1.south west);
    \end{tikzpicture}
    \caption{Ilustração da banda, sombreado e contornado por pontos, e da
    largura de banda, linha contínua, para duas matrizes simétricas.}
    \label{fig:exem_bandwidth}
\end{figure}

O envelope da matriz $A$ é definido como
\begin{align*}
    \text{Env}(A) &= \left\{ \left\{ i, j \right\} \mid 0 < i - j \leq
    \beta_i(A) \right\}.
\end{align*}
e ilustrado na Figura~\ref{fig:exem_profile}.
\begin{figure}[!htb]
    \centering
    \begin{tikzpicture}
        \matrix (A) [matrix of math nodes,%
        left delimiter  = \lbrack,%
        right delimiter = \lbrack] at (0,0)
        {%
            1 & 1 & 1 & 1 & 1 & 0 & 0 \\
            1 & 2 & 1 & 1 & 1 & 1 & 1 \\
            1 & 1 & 2 & 1 & 1 & 1 & 1 \\
            1 & 1 & 1 & 2 & 1 & 0 & 0 \\
            1 & 1 & 1 & 1 & 2 & 0 & 0 \\
            0 & 1 & 1 & 0 & 0 & 3 & 2 \\
            0 & 1 & 1 & 0 & 0 & 2 & 3 \\
        };
        \node[above, shift={(0,.5)}] at (A-1-1) {$1$};
        \node[above, shift={(0,.5)}] at (A-1-2) {$2$};
        \node[above, shift={(0,.5)}] at (A-1-3) {$3$};
        \node[above, shift={(0,.5)}] at (A-1-4) {$4$};
        \node[above, shift={(0,.5)}] at (A-1-5) {$5$};
        \node[above, shift={(0,.5)}] at (A-1-6) {$6$};
        \node[above, shift={(0,.5)}] at (A-1-7) {$7$};
        \node[right, shift={(-1.5,0)}] at (A-1-1) {$1$};
        \node[right, shift={(-1.5,0)}] at (A-2-1) {$2$};
        \node[right, shift={(-1.5,0)}] at (A-3-1) {$3$};
        \node[right, shift={(-1.5,0)}] at (A-4-1) {$4$};
        \node[right, shift={(-1.5,0)}] at (A-5-1) {$5$};
        \node[right, shift={(-1.5,0)}] at (A-6-1) {$6$};
        \node[right, shift={(-1.5,0)}] at (A-7-1) {$7$};
        % Perfil
        \fill[gray, fill opacity=0.2, draw=black, dotted] (A-2-1.north west) -- (A-6-1.north
        west) -- (A-6-2.north west) --
        (A-7-2.south west) -- (A-7-7.south west)
        \foreach \x in {6,5,4,3,2}{
        -- (A-\x-\x.south east) -- (A-\x-\x.south west) -- (A-\x-\x.north west)
        } -- (A-2-1.north west);

        \matrix (B) [matrix of math nodes,%
        left delimiter  = \lbrack,%
        right delimiter = \rbrack] at (8,0)
        {%
           10 &  0 &  2 &  3 &  0 &  0 &  0 \\
            0 & 10 &  2 &  3 &  0 &  0 &  0 \\
            2 &  2 & 10 &  0 &  4 &  0 &  0 \\
            3 &  3 &  0 & 10 &  4 &  0 &  0 \\
            0 &  0 &  4 &  4 & 10 &  2 &  2 \\
            0 &  0 &  0 &  0 &  2 & 10 &  0 \\
            0 &  0 &  0 &  0 &  2 &  0 & 10 \\
        };
        \node[above, shift={(0,.5)}] at (B-1-1) {$1$};
        \node[above, shift={(0,.5)}] at (B-1-2) {$2$};
        \node[above, shift={(0,.5)}] at (B-1-3) {$3$};
        \node[above, shift={(0,.5)}] at (B-1-4) {$4$};
        \node[above, shift={(0,.5)}] at (B-1-5) {$5$};
        \node[above, shift={(0,.5)}] at (B-1-6) {$6$};
        \node[above, shift={(0,.5)}] at (B-1-7) {$7$};
        \node[right, shift={(-1.5,0)}] at (B-1-1) {$1$};
        \node[right, shift={(-1.5,0)}] at (B-2-1) {$2$};
        \node[right, shift={(-1.5,0)}] at (B-3-1) {$3$};
        \node[right, shift={(-1.5,0)}] at (B-4-1) {$4$};
        \node[right, shift={(-1.5,0)}] at (B-5-1) {$5$};
        \node[right, shift={(-1.5,0)}] at (B-6-1) {$6$};
        \node[right, shift={(-1.5,0)}] at (B-7-1) {$7$};
        % Perfil
        \fill[gray, fill opacity=0.2, draw=black, dotted] (B-2-1.south west) -- (B-5-1.north
        west) -- (B-5-3.north west) --
        (B-5-3.south west) -- (B-6-5.north west) -- (B-7-5.south west)
        -- (B-7-7.south west)
        \foreach \x in {6,5,4,3}{
        -- (B-\x-\x.south east) -- (B-\x-\x.south west) -- (B-\x-\x.north west)
        } -- (B-2-1.south west);
    \end{tikzpicture}
    \caption{Ilustração do envelope, sombreado e contornado por
    pontos, para duas matrizes simétricas.}
    \label{fig:exem_profile}
\end{figure}

\subsection{Fatoração de Cholesky}
Seja $A \in \mathbb{R}^{n \times n}$ uma matriz simétrica definida positiva, então existe
uma única matriz $G \in \mathbb{R}^{n \times n}$ triangular inferior com
elementos positivos na diagonal tal que $A = G G^T$. A fatoração $A = G G^T$ é
conhecida como fatoração de Cholesky e no algoritmo a seguir é apresentado
uma forma de obter-se $G$.
\begin{algorithm}[hbt]
    \caption{Pseudo-código da Fatoração de Cholesky}
    \label{alg:chol}
    \begin{algorithmic}[1]
        \REQUIRE $A \in \mathbb{R}^{n \times n}$ simétrica definida positiva.
        \ENSURE $G \in \mathbb{R}^{n \times n}$
        \STATE $G \longleftarrow 0$ 
        \FOR{$k \in 1, \ldots, n$}
            \STATE $G_{k,k} \longleftarrow \sqrt{A_{k,k}}$ \label{alg:chol:sqrt}
            \FOR{$i \in k+1, \ldots n$}
                \STATE $G_{i,k} \longleftarrow A_{i,k} / G_{k,k}$
                \label{alg:chol:div}
            \ENDFOR
            \FOR{$j \in k+1, \ldots, n$}
                \FOR{$i \in k+1, \ldots, n$}
                    \STATE $A_{i,j} \longleftarrow A_{i,j} - G_{i,k} G_{j,k}$
                    \label{alg:chol:prod}
                \ENDFOR
            \ENDFOR
        \ENDFOR
    \end{algorithmic}
\end{algorithm}

\begin{prop}
    \label{prop:pres:env}
    A fatoração de Cholesky preserva o envelope da matriz.
\end{prop}
\begin{proof}
    Para calcular $G$ no Algoritmo~\ref{alg:chol} ocorrem operações aritméticas
    apenas nas linhas \ref{alg:chol:sqrt}, \ref{alg:chol:div} e
    \ref{alg:chol:prod}. Vamos mostrar que nenhuma dessas operações aritméticas
    altera o envelope da matriz.

    Para a linha \ref{alg:chol:sqrt}, como $A$ é simétrica definida positiva,
    ocorre que $A_{k,k}$ é não nulo e, portanto, $G_{k,k}$ também.

    Para a linha \ref{alg:chol:div} é fácil verificar que $G_{i,k}$ é nulo se e
    somente se $A_{i,k}$ também for nulo.

    Para a linha \ref{alg:chol:prod}, se $A_{i,j}$ não pertence ao envelope de $A$
    então $A_{i,j}$ e $G_{i,k}$ são nulos e portanto o envelope não é alterado.
\end{proof}

\begin{prop}
    \label{prop:pres:band}
    A fatoração de Cholesky preserva a banda da matriz.
\end{prop}
\begin{proof}
    Como o envelope da matriz está contido na banda da matriz e já provamos que
    a fatoração de Cholesky preserva o envelope podemos concluir que a fatoração
    de Cholesky também preserva a banda da matriz.
\end{proof}

\section{Método Cuthill-McKee Reverso}
Considere a matriz $A$ e seu fator de Cholesky $G$ dados por
\begin{align*}
    A &= \begin{bmatrix}
        1 &  1 &  1 &  1 &  1 &  1 &  1 \\
        1 & 10 &  0 &  0 &  0 &  0 &  0 \\
        1 &  0 & 10 &  0 &  0 &  0 &  0 \\
        1 &  0 &  0 & 10 &  0 &  0 &  0 \\
        1 &  0 &  0 &  0 & 10 &  0 &  0 \\
        1 &  0 &  0 &  0 &  0 & 10 &  0 \\
        1 &  0 &  0 &  0 &  0 &  0 & 10
    \end{bmatrix}, & G &= \begin{bmatrix}
        1,0 &  0,0 &  0,0 &  0,0 &  0,0 &  0,0 &  0,0 \\
        1,0 &  3,0 &  0,0 &  0,0 &  0,0 &  0,0 &  0,0 \\
        1,0 & -0,3 &  2,9 &  0,0 &  0,0 &  0,0 &  0,0 \\
        1,0 & -0,3 & -0,3 &  2,9 &  0,0 &  0,0 &  0,0 \\
        1,0 & -0,3 & -0,3 & -0,4 &  2,9 &  0,0 &  0,0 \\
        1,0 & -0,3 & -0,3 & -0,4 & -0,4 &  2,8 &  0,0 \\
        1,0 & -0,3 & -0,3 & -0,4 & -0,4 & -0,5 &  2,8
    \end{bmatrix}.
\end{align*}
Observa-se que vários elementos nulos na parte triangular inferior de $A$ são
não nulos em $G$. A perda de elementos nulos ao realizar a fatoração de
Cholesky é denominada de preenchimento e deve ser evitada sempre que possível.

Na seção anterior verificou-se que a fatoração de Cholesky preserva a banda e
envelope da matriz. Por esse motivo, ao utilizar matrizes de banda e
envelope pequenos o preenchimento da matriz também será pequeno. Será que dada
uma matriz $A \in \mathbb{R}^{n \times n}$ simétrica definida positiva existe
uma matriz $A' = P A P^T$, onde $P$ é uma matriz de permutação, tal que a banda
e/ou envelope de $A'$ é menor que o de $A$?

\begin{exem}
    Considere a matriz $A$ e $G$ dada por
    \begin{align*}
        A &= \begin{bmatrix}
            1 & 1 & 1 & 1 & 1 & 0 & 0 \\
            1 & 2 & 1 & 1 & 1 & 1 & 1 \\
            1 & 1 & 2 & 1 & 1 & 1 & 1 \\
            1 & 1 & 1 & 2 & 1 & 0 & 0 \\
            1 & 1 & 1 & 1 & 2 & 0 & 0 \\
            0 & 1 & 1 & 0 & 0 & 3 & 2 \\
            0 & 1 & 1 & 0 & 0 & 2 & 3
        \end{bmatrix},
        & G &= \begin{bmatrix}
            1 & 0 & 0 & 0 & 0 & 0 & 0 \\
            1 & 1 & 0 & 0 & 0 & 0 & 0 \\
            1 & 0 & 1 & 0 & 0 & 0 & 0 \\
            1 & 0 & 0 & 1 & 0 & 0 & 0 \\
            1 & 0 & 0 & 0 & 1 & 0 & 0 \\
            0 & 1 & 1 & 0 & 0 & 1 & 0 \\
            0 & 1 & 1 & 0 & 0 & 0 & 1
        \end{bmatrix}.
    \end{align*}
    Verifica-se que a largura da banda de $A$ e $G$ é $5$ e o envelope é $19$.

    Agora, considere a matriz $A' = P A P^T = G' G'^T$, em que $P$ é uma matriz de
    permutação, dada por
    % O vetor de permutação é
    % [4, 3, 0, 2, 1, 6, 5]
    \begin{align*}
        A' &= \begin{bmatrix}
            2 & 1 & 1 & 1 & 1 & 0 & 0 \\
            1 & 2 & 1 & 1 & 1 & 0 & 0 \\
            1 & 1 & 1 & 1 & 1 & 0 & 0 \\
            1 & 1 & 1 & 2 & 1 & 1 & 1 \\
            1 & 1 & 1 & 1 & 2 & 1 & 1 \\
            0 & 0 & 0 & 1 & 1 & 3 & 2 \\
            0 & 0 & 0 & 1 & 1 & 2 & 3
        \end{bmatrix},
        & G' &= \begin{bmatrix}
             1,4 & 0 & 0 & 0 & 0 & 0 & 0 \\
            0,70 & 1,2 & 0 & 0 & 0 & 0 & 0 \\
            0,70 & 0,40 & 0,57 & 0 & 0 & 0 & 0 \\
            0,70 & 0,40 & 0,57 & 1 & 0 & 0 & 0 \\
            0,70 & 0,40 & 0,57 & 0 & 1 & 0 & 0 \\
            0 & 0 & 0 & 1 & 1 & 1 & 0 \\
            0 & 0 & 0 & 1 & 1 & 0 & 1
        \end{bmatrix}.
    \end{align*}
    Verifica-se que a largura da banda de $A$ e $G$ é $4$ e o envelope é $15$.
\end{exem}

De acordo com o exemplo anterior, verifica-se a possibilidade de preprocessar a matriz
de forma a reduzir a banda e envelope.

Na próxima subseção apresentamos algumas definições da Teoria de Grafos que
serão utilizadas no algoritmo estudado para obter uma permutação que busca
reduzir a largura de banda e envelope.

\subsection{Grafos e matrizes esparsas}
Um grafo é, fundamentalmente, um modo de representar uma relação binária entre
objetos. Para o propósito deste trabalho, considere um grafo $G = (V, E)$ como
um conjunto de vértices $V = \{v_1, v_2, \ldots \}$ e um conjunto de
arestas $E = \{e_1, e_2, \ldots \}$, que são representadas por pares
não ordenados, por exemplo, $e_1 = \{v_1 , v_2\}$.

Assim como um grafo, uma matriz também descreve uma relação binária entre objetos
através de seus elementos não nulos. Uma matriz simétrica $A \in \mathbb{R}^{n
\times n}$ induz um grafo $G(A)$, onde os vértices do grafo correspondem as
dimensões da matriz e a aresta $e = \{i, j\}$ existe se e somente se $A_{ij} \neq 0$.
Na figura abaixo é ilustrado a relação entre uma matriz e um grafo.
\begin{figure}[!hbt]
    \centering
    \begin{tikzpicture}
        \matrix (A) [matrix of math nodes,%
        left delimiter  = \lbrack,%
        right delimiter = \rbrack] at (0,0)
        {%
            1 & 1 & 1 & 1 & 1 & 0 & 0 \\
            1 & 2 & 1 & 1 & 1 & 1 & 1 \\
            1 & 1 & 2 & 1 & 1 & 1 & 1 \\
            1 & 1 & 1 & 2 & 1 & 0 & 0 \\
            1 & 1 & 1 & 1 & 2 & 0 & 0 \\
            0 & 1 & 1 & 0 & 0 & 3 & 2 \\
            0 & 1 & 1 & 0 & 0 & 2 & 3 \\
        };
        %Graph
        \node[draw, circle] (1) at (6,0) {1};
        \node[draw, circle] (2) at (8,1) {2};
        \node[draw, circle] (3) at (8,-1) {3};
        \node[draw, circle] (4) at (4,1) {4};
        \node[draw, circle] (5) at (4,-1) {5};
        \node[draw, circle] (6) at (10,1) {6};
        \node[draw, circle] (7) at (10,-1) {7};
        \draw (1) -- (2);
        \draw (1) -- (3);
        \draw (1) -- (4);
        \draw (1) -- (5);
        \draw (2) -- (3);
        \draw (2) -- (4);
        \draw (2) to[out=190, in=60] (5);
        \draw (2) -- (6);
        \draw (2) -- (7);
        \draw (3) to[out=120, in=-10] (4);
        \draw (3) -- (5);
        \draw (3) -- (6);
        \draw (3) -- (7);
        \draw (4) -- (5);
        \draw (6) -- (7);
    \end{tikzpicture}
    \caption{Ilustração do grafo (a direita) correspondente a uma matriz (a
    esquerda).}
    \label{fig:exem_matrix2graph}
\end{figure}

É importante destacar que permutar, simetricamente, linhas e colunas de uma
matriz corresponde a renumerar os vértices do grafo. Na
Figura~\ref{fig:exem_matrix2graph_perm} é ilustrado uma permutação da matriz
presente na Figura~\ref{fig:exem_matrix2graph} e a renumeração dos vértices do
grafo.
\begin{figure}[!hbt]
    \centering
    \begin{tikzpicture}
        \matrix (A) [matrix of math nodes,%
        left delimiter  = \lbrack,%
        right delimiter = \rbrack] at (0,0)
        {%
            2 & 1 & 1 & 1 & 1 & 0 & 0 \\
            1 & 2 & 1 & 1 & 1 & 0 & 0 \\
            1 & 1 & 1 & 1 & 1 & 0 & 0 \\
            1 & 1 & 1 & 2 & 1 & 1 & 1 \\
            1 & 1 & 1 & 1 & 2 & 1 & 1 \\
            0 & 0 & 0 & 1 & 1 & 3 & 2 \\
            0 & 0 & 0 & 1 & 1 & 2 & 3 \\
        };
        %Graph
        \node[draw, circle] (1) at (6,0) {3};
        \node[draw, circle] (2) at (8,1) {4};
        \node[draw, circle] (3) at (8,-1) {5};
        \node[draw, circle] (4) at (4,1) {2};
        \node[draw, circle] (5) at (4,-1) {1};
        \node[draw, circle] (6) at (10,1) {6};
        \node[draw, circle] (7) at (10,-1) {7};
        \draw (1) -- (2);
        \draw (1) -- (3);
        \draw (1) -- (4);
        \draw (1) -- (5);
        \draw (2) -- (3);
        \draw (2) -- (4);
        \draw (2) to[out=190, in=60] (5);
        \draw (2) -- (6);
        \draw (2) -- (7);
        \draw (3) to[out=120, in=-10] (4);
        \draw (3) -- (5);
        \draw (3) -- (6);
        \draw (3) -- (7);
        \draw (4) -- (5);
        \draw (6) -- (7);
    \end{tikzpicture}
    \caption{Ilustração da permutação simétrica de uma matriz e a renumeração
    dos vértices do grafo correspondente.}
    \label{fig:exem_matrix2graph_perm}
\end{figure}

Um resumo dos principais conceitos é dado a seguir:
\begin{description}
    \item[Vértices adjacentes]  dois vértices $v_1$ e $v_2$ são adjacentes
        quando existe uma aresta entre eles, ou seja, $\{v_1, v_2\} \in E$.
    \item[Incidência] uma aresta incide em um vértice se ele for um dos extremos
        da aresta.
    \item[Grau do vértice] número de arestas incidentes no vértice.
    \item[Caminho] sequência de arestas disjuntas $\left( (x_1, x_2), (x_2,
        x_3), \ldots (x_{k - 1}, x_k) \right)$. O tamanho de um caminho é o
        número de arestas que o compõe.
    \item[Grafo conectado (conexo)] possui pelo menos um caminho entre qualquer
        par de vértices.
    \item[Distância entre vértices] ou $d(v_1, v_2)$, número de arestas que
        formam o menor caminho ligando os vértices $v_1$ e $v_2$.
    \item[Excentricidade] ou $l(v)$, maior distância do vértice $v$ a qualquer outro.
    \item[Diâmetro] maior excentricidade dentre os vértices de um grafo.
    \item[Pseudo-diâmetro] corresponde a uma alta excentricidade, porém não
        necessariamente a maior de todas.
    \item[Vértices periféricos] são vértices cuja excentricidade é igual ao
        diâmetro do grafo.
    \item[Vértices pseudo-periféricos] são vértices que apresentam altas
        excentricidades, mas não necessariamente a maior.
    \item[Estrutura de nível] é uma partição do conjunto $V$ em níveis $L_0,
        L_2, \ldots, L_{l(v)}$ tal que
        \begin{enumerate}
            \item todos os vértices adjacentes aos vértices do nível $L_0$
                encontram-se nos níveis $L_0$ ou $L_1$,
            \item todos os vértices adjacentes aos vértices no nível $L_{l(v)}$
                encontram-se nos níveis $L_{l(v)}$ ou $L_{l(v) - 1}$,
            \item para $0 < i < l(v)$, todos os vértices no nível $L_i$
                encontram-se nos níveis $L_{i - 1}$, $L_i$ ou $L_{i + 1}$.
        \end{enumerate}

        A estrutura de nível com raiz no vértice $v$ é corresponde a
        \begin{enumerate}
            \item $L_0 = \left\{ v \right\}$,
            \item para $i > 0$, $L_i$ é o conjunto de vértices adjacentes aos
                vértices presentes no nível $L_{i - 1}$ e ainda não pertencentes
                a nenhum nível.
        \end{enumerate}

        Na Figura~\ref{fig:graph2level_struct} é ilustrado a estrutura de
        nível, dentre as várias possíveis, para um grafo.
\end{description}
\begin{figure}[htb]
    \centering
    \begin{tikzpicture}
        % graph
        \node[draw, circle] (1) at (2,0) {1};
        \node[draw, circle] (2) at (4,1) {2};
        \node[draw, circle] (3) at (4,-1) {3};
        \node[draw, circle] (4) at (0,1) {4};
        \node[draw, circle] (5) at (0,-1) {5};
        \node[draw, circle] (6) at (6,1) {6};
        \node[draw, circle] (7) at (6,-1) {7};
        \draw (1) -- (2);
        \draw (1) -- (3);
        \draw (1) -- (4);
        \draw (1) -- (5);
        \draw (2) -- (3);
        \draw (2) -- (4);
        \draw (2) to[out=190, in=60] (5);
        \draw (2) -- (6);
        \draw (2) -- (7);
        \draw (3) to[out=120, in=-10] (4);
        \draw (3) -- (5);
        \draw (3) -- (6);
        \draw (3) -- (7);
        \draw (4) -- (5);
        \draw (6) -- (7);
        % level structure
        \node[draw, circle] at (10,2) {6}
            child { node[draw, circle] {2}
                child { node[draw, circle] {1} }
                child { node[draw, circle] {4} }
                child { node[draw, circle] {5} } } 
            child { node[draw, circle] {3} }
            child { node[draw, circle] {7} };
        \path[white, text=black] node at (13,2) {$L_1$}
            child { node {$L_2$}
                child { node {$L_3$} } };
    \end{tikzpicture}
    \caption{Grafo (esquerda) e sua estrutura de nível com raiz no vértice 6
    (direita).}
    \label{fig:graph2level_struct}
\end{figure}
\subsection{Método Cuthill-McKee Reverso}
Cuthill e McKee \cite{Cuthill:1969:ReducingBandwidth} propuseram uma heurística de
reordenação, ver algoritmo abaixo, cujo objetivo principal é reduzir a largura
de banda de uma matriz simétrica $A \in \mathbb{R}^{n \times n}$ cujo grafo
é conexo.
\begin{algorithm}[hbt]
    \caption{Pseudo-código de Cuthill-McKee}
    \label{alg:rcm}
    \begin{algorithmic}[1]
        \REQUIRE Grafo $G(A)$ e um vértice inicial $v$.
        \ENSURE $o$, novo ordenamento dos vértices de $G(A)$.
        \STATE Marca todos os vértices como não visitados.
        \STATE $o \longleftarrow \text{ vetor de zeros}$
        \STATE $i \longleftarrow 1$
        \STATE $f \longleftarrow \text{ fila vazia}$
        \STATE Adicionar $v$ na fila $f$.
        \STATE Marca $v$ como visitado.
        \WHILE{$f$ não for vazia}
            \STATE Desenfileira $f$ em $v$.
            \STATE $o_i \longleftarrow v$
            \STATE $i \longleftarrow i + 1$
            \FORALL{vértice $w$ adjacente a $v$, em ordem crescente de grau,}
                \IF{$w$ ainda não foi visitado}
                    \STATE Adicionar $v$ na fila $f$.
                    \STATE Marca $w$ como visitado.
                \ENDIF
            \ENDFOR
        \ENDWHILE
    \end{algorithmic}
\end{algorithm}

Para o caso de uma matriz $A$ cujo grafo não é conexo, Cuthill e McKee propõe
aplicar a mesma heurística para cada uma das componentes conexas.

George \cite{George:1971:ComputerImplementation} verificou
experimentalmente que ao reverter o ordenamento obtido pelo Algoritmo
Cuthill-McKee, \textit{i.e.}, trocando $o_i$ por $o_{n - i + 1}$ para $i = 1, 2,
\ldots, n$, o novo ordenamento mantém a mesma banda mas diminui o
envelope da matriz. Liu e Sherman \cite{Liu:1976:ComparativeCuthillMcKee}
provaram que essa modificação nunca aumenta o envelope da matriz. Essa versão do
algoritmo é conhecida como Método Cuthill-McKee Reverso.

Uma dos parâmetros de entrada do Método Cuthill-McKee (Reverso) é o vértice inicial sendo
que experimentos computacionais \cite{Cuthill:1969:ReducingBandwidth} sugerem que
vértices pseudo-periféricos são bons candidatos.

Uma heurística para encontrar vértices pseudo-periféricos foi proposto por Alan
George e Joseph W. H. Liu \cite{George:1979:NodeFinder}, ver algoritmo a seguir, e
é uma modificação de um algoritmo anteriormente proposto por Gibbs \textit{et.
al} \cite{Gibbs:1976:ReducingBandwidth} baseado na observação de que
\begin{align*}
    y \in \mathcal{L}_{l(x)}(x) \Longrightarrow l(x) \leq l(y),
\end{align*}
\textit{i.e.}, se $y$ pertence ao nível mais elevado da estrutura de níveis com
raiz em $x$ então a excentricidade de $x$ é menor ou igual a $x$. Essa
observação é verificada com maior facilidade em grafos que são uma árvore, ver
figura abaixo.
\begin{figure}[!hbt]
    \centering
    \begin{tikzpicture}
        \node[draw, circle] at (0,0) {1}
            child { node[draw, circle] {2} }
            child { node[draw, circle] {3}
                child { node[draw, circle] {4} }
                child { node[draw, circle] {5} } };
        \node[draw, circle] at (6,0) {5}
            child { node[draw, circle] {3}
                child { node[draw, circle] {4} }
                child { node[draw, circle] {1} 
                    child { node[draw, circle] {2} } } };
    \end{tikzpicture}
    \caption{Grafo e sua estrutura de níveis com raiz em 1 (a esquerda) e sua
    estrutura de níveis com raiz em 5 (a direita).}
    \label{fig:ilus_obser}
\end{figure}
\begin{algorithm}[H]
    \caption{Pseudo-código para encontrar vértice pseudo-periférico, $x$.}
    \label{alg:ppn}
    \begin{algorithmic}[1]
        \REQUIRE Grafo $G(A)$.
        \ENSURE $x$.
        \STATE $r \longleftarrow \text{Nó arbitrário em }G(A)$
        \STATE Construir estrutura de nível a partir de $r$, $\mathcal{L}(r)$.
        \label{alg:ppn:brle}
        \STATE Escolher um vértice $x$ pertencente ao último nível de
        $\mathcal{L}(r)$.
        \STATE Construir estrutura de nível a partir de $x$, $\mathcal{L}(x)$.
        \IF{$l(x) > l(r)$}
            \STATE $r \longleftarrow x$
            \STATE Retorna para a linha~\ref{alg:ppn:brle}
        \ENDIF
    \end{algorithmic}
\end{algorithm}


% Materiais e Métodos
%     Materiais: citar os equipamentos, reagentes e outros ítens utilizados,
%     informando fabricante ou fornecedor
%     Métodos: descrever os procedimentos detalhados, que possam ser
%     reproduzidos com os materiais e equipamentos descritos
\section{Implementação do Método e Testes Computacionais}
O Método Cuthill-McKee Reverso foi implementado, na linguagem C, pelo autor
deste trabalho como um patch para o PCx
(\url{http://pages.cs.wisc.edu/~swright/PCx/}), e a implementação desenvolvida
encontra-se disponível em \url{https://github.com/r-gaia-cs/PCx.git}.

O PCx é um solver de Programação Linear que utilizar a variante de Mehrotra do
Método Preditor Corretor com a estratégia de correção de Gondzio e o Método do
Mínimo Grau Múltiplo de Liu\cite{George:1981:ComputerSolutionPD} (MMD) para o
reordenamento da matriz $A D^{-1} A^T$.

Testou-se a implementação (tendo como último commit aquele cuja identificação
inicia com 259cf79) com todos os problemas das bibliotecas:
\begin{itemize}
  \item ``Netlib LP'' (\url{http://www.netlib.org/lp/}),
  \item ``Kennington'' (\url{http://www.netlib.org/lp/data/kennington/}),
  \item ``Meszaros'' (\url{http://www.sztaki.hu/~meszaros/public_ftp/lptestset/misc/}),
  \item ``PDS'' (\url{http://plato.asu.edu/ftp/lptestset/pds/}),
  \item ``Rail'' (\url{http://plato.asu.edu/ftp/lptestset/rail/}).
\end{itemize}

Juntando os problemas de todas a bibliotecas, foram utilizados mais de 200
problemas dos quais apenas 4 o RCM apresentou um memor número de elementos não
zeros (3 problemas da Netlib e 1 da Meszaros). Em relação a problemas não
resolvidos, destaca-se os problemas da PDS que o MMD resolveu 5 e o RCM 1
dentre 9 problemas.

Informações adicionas sobre o resultado obtido para cada um dos problemas testados
encontram-se nas Tabelas~\ref{tab:resulken}-\ref{tab:resulrail} sendo que
``R'', ``NNZ'', ``IT'' e ``T''
significam, respectivamente, ``código de retorno'', ``número de elementos não
nulos'', ``número de iterações'' e ``tempo computacional'' (em segundos).
%\begin{table}
%  \centering
%  \caption{Resultados experimentais da biblioteca FCTP.}
%  \label{tab:resulfctp}
%  \begin{tabular}{|l|r|r|r|r|r|r|r|r|}
\hline
\multicolumn{1}{|c|}{Problema} & \multicolumn{4}{|c|}{MMD} &         \multicolumn{4}{|c|}{RCM} \\ \hline
\multicolumn{1}{|c|}{Nome} & \multicolumn{1}{|c|}{R} &
        \multicolumn{1}{|c|}{NNZ} & \multicolumn{1}{|c|}{IT} &
        \multicolumn{1}{|c|}{T} & \multicolumn{1}{|c|}{R} &
        \multicolumn{1}{|c|}{NNZ} & \multicolumn{1}{|c|}{IT} &
        \multicolumn{1}{|c|}{T} \\ \hline
bal8x12 & 3 & 0,0000E+00 & 0 & 0,0000E+00 & 3 & 0,0000E+00 & 0 & 0,0000E+00 \\ \hline
bk4x3 & 3 & 0,0000E+00 & 0 & 0,0000E+00 & 3 & 0,0000E+00 & 0 & 0,0000E+00 \\ \hline
gr4x6 & 3 & 0,0000E+00 & 0 & 0,0000E+00 & 3 & 0,0000E+00 & 0 & 0,0000E+00 \\ \hline
n3700 & 3 & 0,0000E+00 & 0 & 0,0000E+00 & 3 & 0,0000E+00 & 0 & 0,0000E+00 \\ \hline
n3701 & 3 & 0,0000E+00 & 0 & 0,0000E+00 & 3 & 0,0000E+00 & 0 & 0,0000E+00 \\ \hline
n3702 & 3 & 0,0000E+00 & 0 & 0,0000E+00 & 3 & 0,0000E+00 & 0 & 0,0000E+00 \\ \hline
n3703 & 3 & 0,0000E+00 & 0 & 0,0000E+00 & 3 & 0,0000E+00 & 0 & 0,0000E+00 \\ \hline
n3704 & 3 & 0,0000E+00 & 0 & 0,0000E+00 & 3 & 0,0000E+00 & 0 & 0,0000E+00 \\ \hline
n3705 & 3 & 0,0000E+00 & 0 & 0,0000E+00 & 3 & 0,0000E+00 & 0 & 0,0000E+00 \\ \hline
n3706 & 3 & 0,0000E+00 & 0 & 0,0000E+00 & 3 & 0,0000E+00 & 0 & 0,0000E+00 \\ \hline
n3707 & 3 & 0,0000E+00 & 0 & 0,0000E+00 & 3 & 0,0000E+00 & 0 & 0,0000E+00 \\ \hline
n3708 & 3 & 0,0000E+00 & 0 & 0,0000E+00 & 3 & 0,0000E+00 & 0 & 0,0000E+00 \\ \hline
n3709 & 3 & 0,0000E+00 & 0 & 0,0000E+00 & 3 & 0,0000E+00 & 0 & 0,0000E+00 \\ \hline
n370a & 3 & 0,0000E+00 & 0 & 0,0000E+00 & 3 & 0,0000E+00 & 0 & 0,0000E+00 \\ \hline
n370b & 3 & 0,0000E+00 & 0 & 0,0000E+00 & 3 & 0,0000E+00 & 0 & 0,0000E+00 \\ \hline
n370c & 3 & 0,0000E+00 & 0 & 0,0000E+00 & 3 & 0,0000E+00 & 0 & 0,0000E+00 \\ \hline
n370d & 3 & 0,0000E+00 & 0 & 0,0000E+00 & 3 & 0,0000E+00 & 0 & 0,0000E+00 \\ \hline
n370e & 3 & 0,0000E+00 & 0 & 0,0000E+00 & 3 & 0,0000E+00 & 0 & 0,0000E+00 \\ \hline
ran10x10a & 3 & 0,0000E+00 & 0 & 0,0000E+00 & 3 & 0,0000E+00 & 0 & 0,0000E+00 \\ \hline
ran10x10b & 3 & 0,0000E+00 & 0 & 0,0000E+00 & 3 & 0,0000E+00 & 0 & 0,0000E+00 \\ \hline
ran10x10c & 3 & 0,0000E+00 & 0 & 0,0000E+00 & 3 & 0,0000E+00 & 0 & 0,0000E+00 \\ \hline
ran10x12 & 3 & 0,0000E+00 & 0 & 0,0000E+00 & 3 & 0,0000E+00 & 0 & 0,0000E+00 \\ \hline
ran10x26 & 3 & 0,0000E+00 & 0 & 0,0000E+00 & 3 & 0,0000E+00 & 0 & 0,0000E+00 \\ \hline
ran12x12 & 3 & 0,0000E+00 & 0 & 0,0000E+00 & 3 & 0,0000E+00 & 0 & 0,0000E+00 \\ \hline
ran12x21 & 3 & 0,0000E+00 & 0 & 0,0000E+00 & 3 & 0,0000E+00 & 0 & 0,0000E+00 \\ \hline
ran13x13 & 3 & 0,0000E+00 & 0 & 0,0000E+00 & 3 & 0,0000E+00 & 0 & 0,0000E+00 \\ \hline
ran14x18 & 3 & 0,0000E+00 & 0 & 0,0000E+00 & 3 & 0,0000E+00 & 0 & 0,0000E+00 \\ \hline
ran16x16 & 3 & 0,0000E+00 & 0 & 0,0000E+00 & 3 & 0,0000E+00 & 0 & 0,0000E+00 \\ \hline
ran17x17 & 3 & 0,0000E+00 & 0 & 0,0000E+00 & 3 & 0,0000E+00 & 0 & 0,0000E+00 \\ \hline
ran4x64 & 3 & 0,0000E+00 & 0 & 0,0000E+00 & 3 & 0,0000E+00 & 0 & 0,0000E+00 \\ \hline
ran6x43 & 3 & 0,0000E+00 & 0 & 0,0000E+00 & 3 & 0,0000E+00 & 0 & 0,0000E+00 \\ \hline
ran8x32 & 3 & 0,0000E+00 & 0 & 0,0000E+00 & 3 & 0,0000E+00 & 0 & 0,0000E+00 \\ \hline
\end{tabular}

%\end{table}
\begin{table}
  \centering
  \caption{Resultados experimentais da biblioteca Kennington.}
  \label{tab:resulken}
  \begin{tabular}{|l|r|r|r|r|r|r|r|r|}
\hline
\multicolumn{1}{|c|}{Problema} & \multicolumn{4}{|c|}{MMD} &         \multicolumn{4}{|c|}{RCM} \\ \hline
\multicolumn{1}{|c|}{Nome} & \multicolumn{1}{|c|}{R} &
        \multicolumn{1}{|c|}{NNZ} & \multicolumn{1}{|c|}{IT} &
        \multicolumn{1}{|c|}{T} & \multicolumn{1}{|c|}{R} &
        \multicolumn{1}{|c|}{NNZ} & \multicolumn{1}{|c|}{IT} &
        \multicolumn{1}{|c|}{T} \\ \hline
cre-a & 0 & 3,3212E+04 & 23 & 1,8000E-01 & 0 & 2,6381E+06 & 18 & 2,3900E+01 \\ \hline
cre-b & 0 & 2,4863E+05 & 41 & 2,0900E+00 & 0 & 5,3999E+06 & 30 & 1,0192E+02 \\ \hline
cre-c & 0 & 2,8528E+04 & 25 & 1,5000E-01 & 0 & 3,4284E+05 & 23 & 1,9000E+00 \\ \hline
cre-d & 0 & 2,1209E+05 & 41 & 1,7300E+00 & 0 & 5,0801E+06 & 30 & 1,0852E+02 \\ \hline
ken-07 & 0 & 1,0034E+04 & 13 & 4,0000E-02 & 0 & 9,6033E+04 & 13 & 2,8000E-01 \\ \hline
ken-11 & 0 & 1,0291E+05 & 20 & 4,5000E-01 & 0 & 1,0638E+07 & 15 & 2,9063E+02 \\ \hline
ken-13 & 0 & 2,9842E+05 & 23 & 1,2400E+00 & 0 & 2,6536E+07 & 20 & 1,4141E+03 \\ \hline
ken-18 & 0 & 1,9289E+06 & 29 & 8,5100E+00 & -1 & 0,0000E+00 & 0 & 0,0000E+00 \\ \hline
osa-07 & 0 & 2,8276E+04 & 22 & 5,2000E-01 & 0 & 1,8898E+05 & 22 & 1,0100E+00 \\ \hline
osa-14 & 0 & 6,0795E+04 & 25 & 1,3900E+00 & 0 & 8,5696E+05 & 24 & 6,7400E+00 \\ \hline
osa-30 & 0 & 1,1508E+05 & 24 & 2,9800E+00 & 0 & 3,1262E+06 & 19 & 3,5540E+01 \\ \hline
osa-60 & 0 & 2,6591E+05 & 33 & 1,0890E+01 & 0 & 1,5942E+07 & 20 & 4,9131E+02 \\ \hline
pds-02 & 0 & 4,4301E+04 & 24 & 1,9000E-01 & 0 & 7,6833E+05 & 20 & 5,1300E+00 \\ \hline
pds-06 & 0 & 5,8934E+05 & 31 & 4,3900E+00 & 0 & 7,7159E+06 & 27 & 2,4528E+02 \\ \hline
pds-10 & 0 & 1,6877E+06 & 33 & 2,0660E+01 & 0 & 1,5180E+07 & 31 & 6,3432E+02 \\ \hline
pds-20 & 0 & 7,0896E+06 & 43 & 1,8843E+02 & 0 & 4,1662E+07 & 41 & 2,8887E+03 \\ \hline
\end{tabular}

\end{table}
\begin{table}
  \centering
  \caption{Resultados experimentais da biblioteca Meszaros (1/3).}
  \label{tab:resulmes0}
  \begin{tabular}{|l|r|r|r|r|r|r|r|r|}
\hline
\multicolumn{1}{|c|}{Problema} & \multicolumn{4}{|c|}{MMD} &         \multicolumn{4}{|c|}{RCM} \\ \hline
\multicolumn{1}{|c|}{Nome} & \multicolumn{1}{|c|}{R} &
        \multicolumn{1}{|c|}{NNZ} & \multicolumn{1}{|c|}{IT} &
        \multicolumn{1}{|c|}{T} & \multicolumn{1}{|c|}{R} &
        \multicolumn{1}{|c|}{NNZ} & \multicolumn{1}{|c|}{IT} &
        \multicolumn{1}{|c|}{T} \\ \hline
3-mix & 0 & 5,4100E+02 & 16 & 1,0000E-02 & 0 & 6,1300E+02 & 16 & 1,0000E-02 \\ \hline
aa01,mps & 3 & 0,0000E+00 & 0 & 0,0000E+00 & 3 & 0,0000E+00 & 0 & 0,0000E+00 \\ \hline
aa03,mps & 3 & 0,0000E+00 & 0 & 0,0000E+00 & 3 & 0,0000E+00 & 0 & 0,0000E+00 \\ \hline
air02 & 0 & 1,1800E+03 & 18 & 2,7000E-01 & 0 & 1,1900E+03 & 18 & 3,0000E-01 \\ \hline
air03 & 0 & 5,3030E+03 & 34 & 7,1000E-01 & 0 & 5,8610E+03 & 34 & 7,2000E-01 \\ \hline
air04 & 0 & 2,0654E+05 & 76 & 3,7700E+00 & 0 & 2,5434E+05 & 96 & 7,0100E+00 \\ \hline
air05 & 0 & 6,2864E+04 & 33 & 6,1000E-01 & 0 & 7,1846E+04 & 25 & 5,5000E-01 \\ \hline
air06 & 15 & 1,8586E+05 & 53 & 2,3200E+00 & 0 & 2,4812E+05 & 28 & 2,1600E+00 \\ \hline
aircraft & 14 & 3,7540E+03 & 10 & 1,2000E-01 & 14 & 3,7540E+03 & 10 & 1,8000E-01 \\ \hline
aramco & 0 & 9,5000E+01 & 9 & 0,0000E+00 & 0 & 1,1800E+02 & 9 & 0,0000E+00 \\ \hline
bas1lp & 0 & 2,2168E+06 & 9 & 1,8540E+01 & 0 & 3,2825E+06 & 9 & 5,1880E+01 \\ \hline
baxter-mat & 0 & 5,8934E+06 & 30 & 3,9960E+01 & 0 & 2,9029E+07 & 27 & 1,0537E+03 \\ \hline
bpmpd & 6 & 0,0000E+00 & 0 & 0,0000E+00 & 6 & 0,0000E+00 & 0 & 0,0000E+00 \\ \hline
complex1 & 14 & 4,7140E+05 & 16 & 4,0300E+00 & 15 & 4,9191E+05 & 37 & 1,0140E+01 \\ \hline
cr42 & 14 & 1,8040E+03 & 15 & 4,0000E-02 & 14 & 1,8040E+03 & 15 & 4,0000E-02 \\ \hline
crew & 0 & 8,5880E+03 & 12 & 1,6000E-01 & 0 & 8,8230E+03 & 12 & 1,6000E-01 \\ \hline
dano3mip & 12 & 1,0902E+06 & 34 & 1,3600E+01 & 12 & 3,4466E+06 & 32 & 7,8680E+01 \\ \hline
dbah00 & 0 & 2,8099E+05 & 43 & 1,2400E+00 & 0 & 4,8499E+06 & 32 & 1,0051E+02 \\ \hline
dbic1 & 0 & 1,8631E+06 & 45 & 2,8700E+01 & -1 & 0,0000E+00 & 0 & 0,0000E+00 \\ \hline
dbir1 & 0 & 2,5235E+06 & 24 & 6,8810E+01 & 0 & 5,7974E+06 & 23 & 2,8589E+02 \\ \hline
dbir2 & 0 & 2,8185E+06 & 26 & 8,4010E+01 & 0 & 6,7499E+06 & 26 & 3,9022E+02 \\ \hline
delfland & 0 & 3,4610E+03 & 46 & 4,0000E-02 & 0 & 1,2234E+04 & 46 & 8,0000E-02 \\ \hline
disp3,mps & 3 & 0,0000E+00 & 0 & 0,0000E+00 & 3 & 0,0000E+00 & 0 & 0,0000E+00 \\ \hline
ds90 & 0 & 9,9700E+04 & 28 & 3,7000E-01 & 0 & 1,1216E+06 & 21 & 9,9200E+00 \\ \hline
dsbmip & 0 & 1,6789E+04 & 48 & 1,6000E-01 & 0 & 5,9151E+04 & 48 & 4,3000E-01 \\ \hline
emsdz & 0 & 2,7497E+05 & 37 & 1,5000E+00 & 0 & 2,7832E+06 & 29 & 5,1400E+01 \\ \hline
f2177 & 0 & 1,2932E+07 & 2 & 2,8520E+01 & 0 & 2,5060E+07 & 2 & 1,0095E+02 \\ \hline
from-lp-file & -1 & 0,0000E+00 & 0 & 0,0000E+00 & -1 & 0,0000E+00 & 0 & 0,0000E+00 \\ \hline
gamsmod & 0 & 1,6610E+03 & 3 & 1,0000E-02 & 0 & 5,8360E+03 & 3 & 0,0000E+00 \\ \hline
iiasa & 0 & 4,5250E+03 & 15 & 3,0000E-02 & 0 & 4,6540E+03 & 15 & 4,0000E-02 \\ \hline
indata & 0 & 1,0008E+06 & 14 & 5,4600E+00 & 0 & 8,3833E+05 & 14 & 1,0150E+01 \\ \hline
jendrec1 & 6 & 0,0000E+00 & 0 & 0,0000E+00 & 6 & 0,0000E+00 & 0 & 0,0000E+00 \\ \hline
kl02,mps & 3 & 0,0000E+00 & 0 & 0,0000E+00 & 3 & 0,0000E+00 & 0 & 0,0000E+00 \\ \hline
kleemin3 & 0 & 6,0000E+00 & 5 & 0,0000E+00 & 0 & 6,0000E+00 & 5 & 0,0000E+00 \\ \hline
kleemin4 & 0 & 1,0000E+01 & 6 & 0,0000E+00 & 0 & 1,0000E+01 & 6 & 0,0000E+00 \\ \hline
kleemin5 & 0 & 1,5000E+01 & 8 & 0,0000E+00 & 0 & 1,5000E+01 & 8 & 0,0000E+00 \\ \hline
kleemin6 & 0 & 2,1000E+01 & 8 & 0,0000E+00 & 0 & 2,1000E+01 & 8 & 0,0000E+00 \\ \hline
kleemin7 & 0 & 2,8000E+01 & 8 & 0,0000E+00 & 0 & 2,8000E+01 & 8 & 0,0000E+00 \\ \hline
kleemin8 & 0 & 3,6000E+01 & 9 & 0,0000E+00 & 0 & 3,6000E+01 & 9 & 0,0000E+00 \\ \hline
l09a13l1d & 14 & 6,8450E+03 & 12 & 2,0000E-02 & 14 & 7,0910E+03 & 13 & 2,0000E-02 \\ \hline
leader & 0 & 3,9554E+05 & 49 & 2,8600E+00 & 14 & 4,7594E+06 & 48 & 5,1883E+02 \\ \hline
lindo & 0 & 2,0000E+01 & 9 & 0,0000E+00 & 0 & 2,0000E+01 & 9 & 0,0000E+00 \\ \hline
lp22 & 0 & 9,4273E+05 & 56 & 1,5690E+01 & 0 & 2,0149E+06 & 37 & 3,5400E+01 \\ \hline
lpl1 & 0 & 9,6124E+05 & 68 & 1,1760E+01 & 0 & 5,1342E+07 & 52 & 6,7786E+03 \\ \hline
lpl2 & 0 & 4,2869E+04 & 21 & 1,8000E-01 & 0 & 5,0591E+05 & 20 & 3,0700E+00 \\ \hline
lpl3 & 0 & 1,4890E+05 & 36 & 1,1200E+00 & 0 & 4,4995E+06 & 31 & 1,2124E+02 \\ \hline
model10 & 0 & 3,2144E+05 & 45 & 2,0800E+00 & 0 & 1,1938E+06 & 37 & 1,4640E+01 \\ \hline
model11 & 0 & 1,4007E+05 & 25 & 6,4000E-01 & 0 & 7,7321E+05 & 21 & 5,2500E+00 \\ \hline
model1 & 0 & 6,8610E+03 & 8 & 1,0000E-02 & 0 & 1,5093E+04 & 8 & 2,0000E-02 \\ \hline
model2 & 0 & 1,2167E+04 & 23 & 5,0000E-02 & 0 & 1,5749E+04 & 23 & 8,0000E-02 \\ \hline
model3 & 5 & 0,0000E+00 & 0 & 0,0000E+00 & 5 & 0,0000E+00 & 0 & 0,0000E+00 \\ \hline
model4 & 0 & 8,8743E+04 & 35 & 5,0000E-01 & 0 & 1,8571E+05 & 35 & 1,3200E+00 \\ \hline
model5 & 0 & 1,4052E+05 & 44 & 1,2300E+00 & 0 & 2,8807E+05 & 44 & 3,2300E+00 \\ \hline
model6 & 0 & 1,0634E+05 & 33 & 4,3000E-01 & 0 & 3,8909E+05 & 33 & 2,1400E+00 \\ \hline
model7 & 0 & 1,1000E+05 & 39 & 7,1000E-01 & 0 & 2,5601E+05 & 39 & 2,2000E+00 \\ \hline
model8 & 0 & 2,2376E+05 & 17 & 4,2000E-01 & 0 & 2,7580E+05 & 17 & 1,1400E+00 \\ \hline
model9 & 0 & 7,0839E+04 & 41 & 6,0000E-01 & 0 & 1,0504E+05 & 42 & 1,4400E+00 \\ \hline
nemsafm & 0 & 1,0160E+03 & 16 & 1,0000E-02 & 0 & 1,3240E+03 & 16 & 2,0000E-02 \\ \hline
nemscem & 0 & 2,5350E+03 & 18 & 2,0000E-02 & 0 & 1,4749E+04 & 18 & 3,0000E-02 \\ \hline
nemsemm1 & 0 & 2,4298E+05 & 57 & 1,5710E+01 & 0 & 3,9867E+05 & 57 & 1,9620E+01 \\ \hline
nemsemm2 & 0 & 6,2496E+04 & 35 & 8,7000E-01 & 0 & 3,8140E+05 & 35 & 3,3700E+00 \\ \hline
nemspmm1 & 0 & 1,1464E+05 & 39 & 7,4000E-01 & 0 & 8,6560E+05 & 31 & 9,0400E+00 \\ \hline
nemspmm2,in & 0 & 1,1890E+05 & 44 & 1,0100E+00 & 0 & 7,1490E+05 & 34 & 7,4900E+00 \\ \hline
nemswrld & 0 & 7,8312E+05 & 42 & 7,8000E+00 & 0 & 4,4720E+06 & 38 & 9,3030E+01 \\ \hline
nsct1 & 0 & 5,4180E+06 & 21 & 1,7080E+02 & 0 & 8,5935E+06 & 22 & 6,1142E+02 \\ \hline
nsct2 & 0 & 5,6456E+06 & 28 & 2,3305E+02 & 0 & 8,3502E+06 & 30 & 7,0599E+02 \\ \hline
nsic1 & 0 & 5,0420E+03 & 10 & 2,0000E-02 & 0 & 7,2400E+03 & 10 & 3,0000E-02 \\ \hline
nsic2 & 0 & 5,5160E+03 & 14 & 2,0000E-02 & 0 & 8,0110E+03 & 14 & 4,0000E-02 \\ \hline
nsir1 & 0 & 3,0346E+05 & 21 & 3,6600E+00 & 0 & 6,3039E+05 & 17 & 1,0390E+01 \\ \hline
nsir2 & 0 & 3,2795E+05 & 30 & 5,9400E+00 & 0 & 6,5743E+05 & 26 & 1,4420E+01 \\ \hline
nw14,mps & 3 & 0,0000E+00 & 0 & 0,0000E+00 & 3 & 0,0000E+00 & 0 & 0,0000E+00 \\ \hline
olivier & 14 & 4,6137E+05 & 50 & 4,5700E+00 & 14 & 5,6194E+06 & 43 & 4,4814E+02 \\ \hline
or4,mps & 3 & 0,0000E+00 & 0 & 0,0000E+00 & 3 & 0,0000E+00 & 0 & 0,0000E+00 \\ \hline
orna1,mps & 3 & 0,0000E+00 & 0 & 0,0000E+00 & 3 & 0,0000E+00 & 0 & 0,0000E+00 \\ \hline
orna2,mps & 3 & 0,0000E+00 & 0 & 0,0000E+00 & 3 & 0,0000E+00 & 0 & 0,0000E+00 \\ \hline
orna3,mps & 3 & 0,0000E+00 & 0 & 0,0000E+00 & 3 & 0,0000E+00 & 0 & 0,0000E+00 \\ \hline
orna7,mps & 3 & 0,0000E+00 & 0 & 0,0000E+00 & 3 & 0,0000E+00 & 0 & 0,0000E+00 \\ \hline
orswq2,mps & 3 & 0,0000E+00 & 0 & 0,0000E+00 & 3 & 0,0000E+00 & 0 & 0,0000E+00 \\ \hline
p0033 & 0 & 5,8000E+01 & 8 & 0,0000E+00 & 0 & 6,0000E+01 & 8 & 0,0000E+00 \\ \hline
p0040 & 0 & 9,6000E+01 & 9 & 0,0000E+00 & 0 & 9,6000E+01 & 9 & 0,0000E+00 \\ \hline
p0201 & 0 & 2,9220E+03 & 5 & 1,0000E-02 & 0 & 7,8160E+03 & 5 & 1,0000E-02 \\ \hline
p0282 & 0 & 9,5870E+03 & 13 & 2,0000E-02 & 0 & 1,3065E+04 & 13 & 4,0000E-02 \\ \hline
p0291 & 0 & 4,8890E+03 & 12 & 1,0000E-02 & 0 & 8,7050E+03 & 12 & 2,0000E-02 \\ \hline
p0548 & 0 & 1,0530E+03 & 24 & 2,0000E-02 & 0 & 1,1890E+03 & 24 & 2,0000E-02 \\ \hline
p05 & 0 & 2,3647E+05 & 23 & 7,1000E-01 & 0 & 2,2887E+06 & 19 & 2,3310E+01 \\ \hline
p10 & 0 & 4,6962E+05 & 27 & 1,6900E+00 & 0 & 7,7808E+06 & 29 & 2,5367E+02 \\ \hline
p12345 & 0 & 6,4360E+03 & 14 & 8,0000E-02 & 0 & 6,4360E+03 & 14 & 8,0000E-02 \\ \hline
p19328 & 0 & 1,6932E+04 & 15 & 2,7000E-01 & 0 & 1,6932E+04 & 15 & 3,1000E-01 \\ \hline
p19 & 0 & 2,9710E+03 & 19 & 2,0000E-02 & 0 & 9,8150E+03 & 19 & 3,0000E-02 \\ \hline
p2756 & 0 & 6,5240E+03 & 15 & 5,0000E-02 & 0 & 6,7656E+04 & 15 & 2,0000E-01 \\ \hline
pcb1000 & 0 & 2,9753E+04 & 23 & 1,8000E-01 & 0 & 3,9440E+05 & 23 & 1,8200E+00 \\ \hline
pcb3000 & 0 & 1,0109E+05 & 25 & 5,6000E-01 & 0 & 4,3551E+06 & 19 & 4,9320E+01 \\ \hline
primagaz & 0 & 1,2404E+04 & 13 & 1,1000E-01 & 0 & 1,2404E+04 & 13 & 1,4000E-01 \\ \hline
progas & 0 & 3,0946E+04 & 15 & 6,0000E-02 & 0 & 4,8144E+04 & 15 & 2,3000E-01 \\ \hline
ps & 15 & 1,1054E+07 & 40 & 5,1600E+02 & 15 & 1,2683E+07 & 48 & 6,2349E+02 \\ \hline
qiu & 0 & 4,6492E+04 & 5 & 2,0000E-02 & 0 & 1,0207E+05 & 5 & 8,0000E-02 \\ \hline
r05 & 0 & 4,3308E+05 & 23 & 1,7800E+00 & 0 & 2,1250E+06 & 18 & 2,6230E+01 \\ \hline
radio & -1 & 0,0000E+00 & 0 & 0,0000E+00 & -11 & 0,0000E+00 & 0 & 0,0000E+00 \\ \hline
rlf,pre,dual & 0 & 8,4078E+05 & 13 & 2,3400E+00 & 0 & 6,0951E+06 & 13 & 4,8841E+02 \\ \hline
routing & 0 & 3,1143E+06 & 20 & 6,0600E+00 & 0 & 1,2819E+07 & 17 & 8,7570E+01 \\ \hline
sc205 & 12 & 5,0422E+04 & 15 & 5,3000E-01 & 12 & 5,3625E+04 & 15 & 6,0000E-01 \\ \hline
seymour & 0 & 4,5141E+06 & 15 & 3,6330E+01 & 0 & 8,2165E+06 & 15 & 1,3423E+02 \\ \hline
slp-tsk & 0 & 4,0878E+06 & 19 & 6,4860E+01 & 0 & 4,0918E+06 & 19 & 1,2514E+02 \\ \hline
southern1 & 0 & 5,6710E+06 & 15 & 6,5000E+01 & 0 & 7,7658E+06 & 15 & 6,6628E+02 \\ \hline
sturing & 0 & 1,8861E+04 & 36 & 1,5000E-01 & 0 & 3,6899E+04 & 36 & 3,0000E-01 \\ \hline
sws & 0 & 5,8235E+04 & 10 & 2,1000E-01 & 0 & 4,6605E+05 & 10 & 5,9000E-01 \\ \hline
t0331-4l & 0 & 2,0167E+05 & 37 & 6,5800E+00 & 0 & 2,1556E+05 & 37 & 7,0600E+00 \\ \hline
test & 0 & 6,5540E+03 & 46 & 2,7000E-01 & 0 & 6,5540E+03 & 46 & 2,9000E-01 \\ \hline
testps,mod & 0 & 3,1000E+01 & 13 & 0,0000E+00 & 0 & 3,1000E+01 & 13 & 0,0000E+00 \\ \hline
unilever2 & 15 & 4,2276E+05 & 36 & 3,4700E+00 & 0 & 5,8050E+06 & 73 & 2,4664E+02 \\ \hline
us04,mps & 3 & 0,0000E+00 & 0 & 0,0000E+00 & 3 & 0,0000E+00 & 0 & 0,0000E+00 \\ \hline
world & 14 & 1,1112E+06 & 60 & 9,2800E+00 & -1 & 0,0000E+00 & 0 & 0,0000E+00 \\ \hline
zed & 0 & 3,3800E+03 & 9 & 2,0000E-02 & 0 & 3,9520E+03 & 9 & 2,0000E-02 \\ \hline
zz & 15 & 2,8570E+05 & 43 & 3,7000E+00 & 15 & 9,9260E+05 & 43 & 1,0700E+01 \\ \hline
\end{tabular}

\end{table}
\begin{table}
  \centering
  \caption{Resultados experimentais da biblioteca Meszaros (2/3).}
  \label{tab:resulmes1}
  \begin{tabular}{|l|r|r|r|r|r|r|r|r|}
\hline
\multicolumn{1}{|c|}{Problema} & \multicolumn{4}{|c|}{MMD} &         \multicolumn{4}{|c|}{RCM} \\ \hline
\multicolumn{1}{|c|}{Nome} & \multicolumn{1}{|c|}{R} &
        \multicolumn{1}{|c|}{NNZ} & \multicolumn{1}{|c|}{IT} &
        \multicolumn{1}{|c|}{T} & \multicolumn{1}{|c|}{R} &
        \multicolumn{1}{|c|}{NNZ} & \multicolumn{1}{|c|}{IT} &
        \multicolumn{1}{|c|}{T} \\ \hline
3-mix & 0 & 5,4100E+02 & 16 & 1,0000E-02 & 0 & 6,1300E+02 & 16 & 1,0000E-02 \\ \hline
aa01,mps & 3 & 0,0000E+00 & 0 & 0,0000E+00 & 3 & 0,0000E+00 & 0 & 0,0000E+00 \\ \hline
aa03,mps & 3 & 0,0000E+00 & 0 & 0,0000E+00 & 3 & 0,0000E+00 & 0 & 0,0000E+00 \\ \hline
air02 & 0 & 1,1800E+03 & 18 & 2,7000E-01 & 0 & 1,1900E+03 & 18 & 3,0000E-01 \\ \hline
air03 & 0 & 5,3030E+03 & 34 & 7,1000E-01 & 0 & 5,8610E+03 & 34 & 7,2000E-01 \\ \hline
air04 & 0 & 2,0654E+05 & 76 & 3,7700E+00 & 0 & 2,5434E+05 & 96 & 7,0100E+00 \\ \hline
air05 & 0 & 6,2864E+04 & 33 & 6,1000E-01 & 0 & 7,1846E+04 & 25 & 5,5000E-01 \\ \hline
air06 & 15 & 1,8586E+05 & 53 & 2,3200E+00 & 0 & 2,4812E+05 & 28 & 2,1600E+00 \\ \hline
aircraft & 14 & 3,7540E+03 & 10 & 1,2000E-01 & 14 & 3,7540E+03 & 10 & 1,8000E-01 \\ \hline
aramco & 0 & 9,5000E+01 & 9 & 0,0000E+00 & 0 & 1,1800E+02 & 9 & 0,0000E+00 \\ \hline
bas1lp & 0 & 2,2168E+06 & 9 & 1,8540E+01 & 0 & 3,2825E+06 & 9 & 5,1880E+01 \\ \hline
baxter-mat & 0 & 5,8934E+06 & 30 & 3,9960E+01 & 0 & 2,9029E+07 & 27 & 1,0537E+03 \\ \hline
bpmpd & 6 & 0,0000E+00 & 0 & 0,0000E+00 & 6 & 0,0000E+00 & 0 & 0,0000E+00 \\ \hline
complex1 & 14 & 4,7140E+05 & 16 & 4,0300E+00 & 15 & 4,9191E+05 & 37 & 1,0140E+01 \\ \hline
cr42 & 14 & 1,8040E+03 & 15 & 4,0000E-02 & 14 & 1,8040E+03 & 15 & 4,0000E-02 \\ \hline
crew & 0 & 8,5880E+03 & 12 & 1,6000E-01 & 0 & 8,8230E+03 & 12 & 1,6000E-01 \\ \hline
dano3mip & 12 & 1,0902E+06 & 34 & 1,3600E+01 & 12 & 3,4466E+06 & 32 & 7,8680E+01 \\ \hline
dbah00 & 0 & 2,8099E+05 & 43 & 1,2400E+00 & 0 & 4,8499E+06 & 32 & 1,0051E+02 \\ \hline
dbic1 & 0 & 1,8631E+06 & 45 & 2,8700E+01 & -1 & 0,0000E+00 & 0 & 0,0000E+00 \\ \hline
dbir1 & 0 & 2,5235E+06 & 24 & 6,8810E+01 & 0 & 5,7974E+06 & 23 & 2,8589E+02 \\ \hline
dbir2 & 0 & 2,8185E+06 & 26 & 8,4010E+01 & 0 & 6,7499E+06 & 26 & 3,9022E+02 \\ \hline
delfland & 0 & 3,4610E+03 & 46 & 4,0000E-02 & 0 & 1,2234E+04 & 46 & 8,0000E-02 \\ \hline
disp3,mps & 3 & 0,0000E+00 & 0 & 0,0000E+00 & 3 & 0,0000E+00 & 0 & 0,0000E+00 \\ \hline
ds90 & 0 & 9,9700E+04 & 28 & 3,7000E-01 & 0 & 1,1216E+06 & 21 & 9,9200E+00 \\ \hline
dsbmip & 0 & 1,6789E+04 & 48 & 1,6000E-01 & 0 & 5,9151E+04 & 48 & 4,3000E-01 \\ \hline
emsdz & 0 & 2,7497E+05 & 37 & 1,5000E+00 & 0 & 2,7832E+06 & 29 & 5,1400E+01 \\ \hline
f2177 & 0 & 1,2932E+07 & 2 & 2,8520E+01 & 0 & 2,5060E+07 & 2 & 1,0095E+02 \\ \hline
from-lp-file & -1 & 0,0000E+00 & 0 & 0,0000E+00 & -1 & 0,0000E+00 & 0 & 0,0000E+00 \\ \hline
gamsmod & 0 & 1,6610E+03 & 3 & 1,0000E-02 & 0 & 5,8360E+03 & 3 & 0,0000E+00 \\ \hline
iiasa & 0 & 4,5250E+03 & 15 & 3,0000E-02 & 0 & 4,6540E+03 & 15 & 4,0000E-02 \\ \hline
indata & 0 & 1,0008E+06 & 14 & 5,4600E+00 & 0 & 8,3833E+05 & 14 & 1,0150E+01 \\ \hline
jendrec1 & 6 & 0,0000E+00 & 0 & 0,0000E+00 & 6 & 0,0000E+00 & 0 & 0,0000E+00 \\ \hline
kl02,mps & 3 & 0,0000E+00 & 0 & 0,0000E+00 & 3 & 0,0000E+00 & 0 & 0,0000E+00 \\ \hline
kleemin3 & 0 & 6,0000E+00 & 5 & 0,0000E+00 & 0 & 6,0000E+00 & 5 & 0,0000E+00 \\ \hline
kleemin4 & 0 & 1,0000E+01 & 6 & 0,0000E+00 & 0 & 1,0000E+01 & 6 & 0,0000E+00 \\ \hline
kleemin5 & 0 & 1,5000E+01 & 8 & 0,0000E+00 & 0 & 1,5000E+01 & 8 & 0,0000E+00 \\ \hline
kleemin6 & 0 & 2,1000E+01 & 8 & 0,0000E+00 & 0 & 2,1000E+01 & 8 & 0,0000E+00 \\ \hline
kleemin7 & 0 & 2,8000E+01 & 8 & 0,0000E+00 & 0 & 2,8000E+01 & 8 & 0,0000E+00 \\ \hline
kleemin8 & 0 & 3,6000E+01 & 9 & 0,0000E+00 & 0 & 3,6000E+01 & 9 & 0,0000E+00 \\ \hline
l09a13l1d & 14 & 6,8450E+03 & 12 & 2,0000E-02 & 14 & 7,0910E+03 & 13 & 2,0000E-02 \\ \hline
leader & 0 & 3,9554E+05 & 49 & 2,8600E+00 & 14 & 4,7594E+06 & 48 & 5,1883E+02 \\ \hline
lindo & 0 & 2,0000E+01 & 9 & 0,0000E+00 & 0 & 2,0000E+01 & 9 & 0,0000E+00 \\ \hline
lp22 & 0 & 9,4273E+05 & 56 & 1,5690E+01 & 0 & 2,0149E+06 & 37 & 3,5400E+01 \\ \hline
lpl1 & 0 & 9,6124E+05 & 68 & 1,1760E+01 & 0 & 5,1342E+07 & 52 & 6,7786E+03 \\ \hline
lpl2 & 0 & 4,2869E+04 & 21 & 1,8000E-01 & 0 & 5,0591E+05 & 20 & 3,0700E+00 \\ \hline
lpl3 & 0 & 1,4890E+05 & 36 & 1,1200E+00 & 0 & 4,4995E+06 & 31 & 1,2124E+02 \\ \hline
model10 & 0 & 3,2144E+05 & 45 & 2,0800E+00 & 0 & 1,1938E+06 & 37 & 1,4640E+01 \\ \hline
model11 & 0 & 1,4007E+05 & 25 & 6,4000E-01 & 0 & 7,7321E+05 & 21 & 5,2500E+00 \\ \hline
model1 & 0 & 6,8610E+03 & 8 & 1,0000E-02 & 0 & 1,5093E+04 & 8 & 2,0000E-02 \\ \hline
model2 & 0 & 1,2167E+04 & 23 & 5,0000E-02 & 0 & 1,5749E+04 & 23 & 8,0000E-02 \\ \hline
model3 & 5 & 0,0000E+00 & 0 & 0,0000E+00 & 5 & 0,0000E+00 & 0 & 0,0000E+00 \\ \hline
model4 & 0 & 8,8743E+04 & 35 & 5,0000E-01 & 0 & 1,8571E+05 & 35 & 1,3200E+00 \\ \hline
model5 & 0 & 1,4052E+05 & 44 & 1,2300E+00 & 0 & 2,8807E+05 & 44 & 3,2300E+00 \\ \hline
model6 & 0 & 1,0634E+05 & 33 & 4,3000E-01 & 0 & 3,8909E+05 & 33 & 2,1400E+00 \\ \hline
model7 & 0 & 1,1000E+05 & 39 & 7,1000E-01 & 0 & 2,5601E+05 & 39 & 2,2000E+00 \\ \hline
model8 & 0 & 2,2376E+05 & 17 & 4,2000E-01 & 0 & 2,7580E+05 & 17 & 1,1400E+00 \\ \hline
model9 & 0 & 7,0839E+04 & 41 & 6,0000E-01 & 0 & 1,0504E+05 & 42 & 1,4400E+00 \\ \hline
nemsafm & 0 & 1,0160E+03 & 16 & 1,0000E-02 & 0 & 1,3240E+03 & 16 & 2,0000E-02 \\ \hline
nemscem & 0 & 2,5350E+03 & 18 & 2,0000E-02 & 0 & 1,4749E+04 & 18 & 3,0000E-02 \\ \hline
nemsemm1 & 0 & 2,4298E+05 & 57 & 1,5710E+01 & 0 & 3,9867E+05 & 57 & 1,9620E+01 \\ \hline
nemsemm2 & 0 & 6,2496E+04 & 35 & 8,7000E-01 & 0 & 3,8140E+05 & 35 & 3,3700E+00 \\ \hline
nemspmm1 & 0 & 1,1464E+05 & 39 & 7,4000E-01 & 0 & 8,6560E+05 & 31 & 9,0400E+00 \\ \hline
nemspmm2,in & 0 & 1,1890E+05 & 44 & 1,0100E+00 & 0 & 7,1490E+05 & 34 & 7,4900E+00 \\ \hline
nemswrld & 0 & 7,8312E+05 & 42 & 7,8000E+00 & 0 & 4,4720E+06 & 38 & 9,3030E+01 \\ \hline
nsct1 & 0 & 5,4180E+06 & 21 & 1,7080E+02 & 0 & 8,5935E+06 & 22 & 6,1142E+02 \\ \hline
nsct2 & 0 & 5,6456E+06 & 28 & 2,3305E+02 & 0 & 8,3502E+06 & 30 & 7,0599E+02 \\ \hline
nsic1 & 0 & 5,0420E+03 & 10 & 2,0000E-02 & 0 & 7,2400E+03 & 10 & 3,0000E-02 \\ \hline
nsic2 & 0 & 5,5160E+03 & 14 & 2,0000E-02 & 0 & 8,0110E+03 & 14 & 4,0000E-02 \\ \hline
nsir1 & 0 & 3,0346E+05 & 21 & 3,6600E+00 & 0 & 6,3039E+05 & 17 & 1,0390E+01 \\ \hline
nsir2 & 0 & 3,2795E+05 & 30 & 5,9400E+00 & 0 & 6,5743E+05 & 26 & 1,4420E+01 \\ \hline
nw14,mps & 3 & 0,0000E+00 & 0 & 0,0000E+00 & 3 & 0,0000E+00 & 0 & 0,0000E+00 \\ \hline
olivier & 14 & 4,6137E+05 & 50 & 4,5700E+00 & 14 & 5,6194E+06 & 43 & 4,4814E+02 \\ \hline
or4,mps & 3 & 0,0000E+00 & 0 & 0,0000E+00 & 3 & 0,0000E+00 & 0 & 0,0000E+00 \\ \hline
orna1,mps & 3 & 0,0000E+00 & 0 & 0,0000E+00 & 3 & 0,0000E+00 & 0 & 0,0000E+00 \\ \hline
orna2,mps & 3 & 0,0000E+00 & 0 & 0,0000E+00 & 3 & 0,0000E+00 & 0 & 0,0000E+00 \\ \hline
orna3,mps & 3 & 0,0000E+00 & 0 & 0,0000E+00 & 3 & 0,0000E+00 & 0 & 0,0000E+00 \\ \hline
orna7,mps & 3 & 0,0000E+00 & 0 & 0,0000E+00 & 3 & 0,0000E+00 & 0 & 0,0000E+00 \\ \hline
orswq2,mps & 3 & 0,0000E+00 & 0 & 0,0000E+00 & 3 & 0,0000E+00 & 0 & 0,0000E+00 \\ \hline
p0033 & 0 & 5,8000E+01 & 8 & 0,0000E+00 & 0 & 6,0000E+01 & 8 & 0,0000E+00 \\ \hline
p0040 & 0 & 9,6000E+01 & 9 & 0,0000E+00 & 0 & 9,6000E+01 & 9 & 0,0000E+00 \\ \hline
p0201 & 0 & 2,9220E+03 & 5 & 1,0000E-02 & 0 & 7,8160E+03 & 5 & 1,0000E-02 \\ \hline
p0282 & 0 & 9,5870E+03 & 13 & 2,0000E-02 & 0 & 1,3065E+04 & 13 & 4,0000E-02 \\ \hline
p0291 & 0 & 4,8890E+03 & 12 & 1,0000E-02 & 0 & 8,7050E+03 & 12 & 2,0000E-02 \\ \hline
p0548 & 0 & 1,0530E+03 & 24 & 2,0000E-02 & 0 & 1,1890E+03 & 24 & 2,0000E-02 \\ \hline
p05 & 0 & 2,3647E+05 & 23 & 7,1000E-01 & 0 & 2,2887E+06 & 19 & 2,3310E+01 \\ \hline
p10 & 0 & 4,6962E+05 & 27 & 1,6900E+00 & 0 & 7,7808E+06 & 29 & 2,5367E+02 \\ \hline
p12345 & 0 & 6,4360E+03 & 14 & 8,0000E-02 & 0 & 6,4360E+03 & 14 & 8,0000E-02 \\ \hline
p19328 & 0 & 1,6932E+04 & 15 & 2,7000E-01 & 0 & 1,6932E+04 & 15 & 3,1000E-01 \\ \hline
p19 & 0 & 2,9710E+03 & 19 & 2,0000E-02 & 0 & 9,8150E+03 & 19 & 3,0000E-02 \\ \hline
p2756 & 0 & 6,5240E+03 & 15 & 5,0000E-02 & 0 & 6,7656E+04 & 15 & 2,0000E-01 \\ \hline
pcb1000 & 0 & 2,9753E+04 & 23 & 1,8000E-01 & 0 & 3,9440E+05 & 23 & 1,8200E+00 \\ \hline
pcb3000 & 0 & 1,0109E+05 & 25 & 5,6000E-01 & 0 & 4,3551E+06 & 19 & 4,9320E+01 \\ \hline
primagaz & 0 & 1,2404E+04 & 13 & 1,1000E-01 & 0 & 1,2404E+04 & 13 & 1,4000E-01 \\ \hline
progas & 0 & 3,0946E+04 & 15 & 6,0000E-02 & 0 & 4,8144E+04 & 15 & 2,3000E-01 \\ \hline
ps & 15 & 1,1054E+07 & 40 & 5,1600E+02 & 15 & 1,2683E+07 & 48 & 6,2349E+02 \\ \hline
qiu & 0 & 4,6492E+04 & 5 & 2,0000E-02 & 0 & 1,0207E+05 & 5 & 8,0000E-02 \\ \hline
r05 & 0 & 4,3308E+05 & 23 & 1,7800E+00 & 0 & 2,1250E+06 & 18 & 2,6230E+01 \\ \hline
radio & -1 & 0,0000E+00 & 0 & 0,0000E+00 & -11 & 0,0000E+00 & 0 & 0,0000E+00 \\ \hline
rlf,pre,dual & 0 & 8,4078E+05 & 13 & 2,3400E+00 & 0 & 6,0951E+06 & 13 & 4,8841E+02 \\ \hline
routing & 0 & 3,1143E+06 & 20 & 6,0600E+00 & 0 & 1,2819E+07 & 17 & 8,7570E+01 \\ \hline
sc205 & 12 & 5,0422E+04 & 15 & 5,3000E-01 & 12 & 5,3625E+04 & 15 & 6,0000E-01 \\ \hline
seymour & 0 & 4,5141E+06 & 15 & 3,6330E+01 & 0 & 8,2165E+06 & 15 & 1,3423E+02 \\ \hline
slp-tsk & 0 & 4,0878E+06 & 19 & 6,4860E+01 & 0 & 4,0918E+06 & 19 & 1,2514E+02 \\ \hline
southern1 & 0 & 5,6710E+06 & 15 & 6,5000E+01 & 0 & 7,7658E+06 & 15 & 6,6628E+02 \\ \hline
sturing & 0 & 1,8861E+04 & 36 & 1,5000E-01 & 0 & 3,6899E+04 & 36 & 3,0000E-01 \\ \hline
sws & 0 & 5,8235E+04 & 10 & 2,1000E-01 & 0 & 4,6605E+05 & 10 & 5,9000E-01 \\ \hline
t0331-4l & 0 & 2,0167E+05 & 37 & 6,5800E+00 & 0 & 2,1556E+05 & 37 & 7,0600E+00 \\ \hline
test & 0 & 6,5540E+03 & 46 & 2,7000E-01 & 0 & 6,5540E+03 & 46 & 2,9000E-01 \\ \hline
testps,mod & 0 & 3,1000E+01 & 13 & 0,0000E+00 & 0 & 3,1000E+01 & 13 & 0,0000E+00 \\ \hline
unilever2 & 15 & 4,2276E+05 & 36 & 3,4700E+00 & 0 & 5,8050E+06 & 73 & 2,4664E+02 \\ \hline
us04,mps & 3 & 0,0000E+00 & 0 & 0,0000E+00 & 3 & 0,0000E+00 & 0 & 0,0000E+00 \\ \hline
world & 14 & 1,1112E+06 & 60 & 9,2800E+00 & -1 & 0,0000E+00 & 0 & 0,0000E+00 \\ \hline
zed & 0 & 3,3800E+03 & 9 & 2,0000E-02 & 0 & 3,9520E+03 & 9 & 2,0000E-02 \\ \hline
zz & 15 & 2,8570E+05 & 43 & 3,7000E+00 & 15 & 9,9260E+05 & 43 & 1,0700E+01 \\ \hline
\end{tabular}

\end{table}
\begin{table}
  \centering
  \caption{Resultados experimentais da biblioteca Meszaros (3/3).}
  \label{tab:resulmes2}
  \begin{tabular}{|l|r|r|r|r|r|r|r|r|}
\hline
\multicolumn{1}{|c|}{Problema} & \multicolumn{4}{|c|}{MMD} &         \multicolumn{4}{|c|}{RCM} \\ \hline
\multicolumn{1}{|c|}{Nome} & \multicolumn{1}{|c|}{R} &
        \multicolumn{1}{|c|}{NNZ} & \multicolumn{1}{|c|}{IT} &
        \multicolumn{1}{|c|}{T} & \multicolumn{1}{|c|}{R} &
        \multicolumn{1}{|c|}{NNZ} & \multicolumn{1}{|c|}{IT} &
        \multicolumn{1}{|c|}{T} \\ \hline
3-mix & 0 & 5,4100E+02 & 16 & 1,0000E-02 & 0 & 6,1300E+02 & 16 & 1,0000E-02 \\ \hline
aa01,mps & 3 & 0,0000E+00 & 0 & 0,0000E+00 & 3 & 0,0000E+00 & 0 & 0,0000E+00 \\ \hline
aa03,mps & 3 & 0,0000E+00 & 0 & 0,0000E+00 & 3 & 0,0000E+00 & 0 & 0,0000E+00 \\ \hline
air02 & 0 & 1,1800E+03 & 18 & 2,7000E-01 & 0 & 1,1900E+03 & 18 & 3,0000E-01 \\ \hline
air03 & 0 & 5,3030E+03 & 34 & 7,1000E-01 & 0 & 5,8610E+03 & 34 & 7,2000E-01 \\ \hline
air04 & 0 & 2,0654E+05 & 76 & 3,7700E+00 & 0 & 2,5434E+05 & 96 & 7,0100E+00 \\ \hline
air05 & 0 & 6,2864E+04 & 33 & 6,1000E-01 & 0 & 7,1846E+04 & 25 & 5,5000E-01 \\ \hline
air06 & 15 & 1,8586E+05 & 53 & 2,3200E+00 & 0 & 2,4812E+05 & 28 & 2,1600E+00 \\ \hline
aircraft & 14 & 3,7540E+03 & 10 & 1,2000E-01 & 14 & 3,7540E+03 & 10 & 1,8000E-01 \\ \hline
aramco & 0 & 9,5000E+01 & 9 & 0,0000E+00 & 0 & 1,1800E+02 & 9 & 0,0000E+00 \\ \hline
bas1lp & 0 & 2,2168E+06 & 9 & 1,8540E+01 & 0 & 3,2825E+06 & 9 & 5,1880E+01 \\ \hline
baxter-mat & 0 & 5,8934E+06 & 30 & 3,9960E+01 & 0 & 2,9029E+07 & 27 & 1,0537E+03 \\ \hline
bpmpd & 6 & 0,0000E+00 & 0 & 0,0000E+00 & 6 & 0,0000E+00 & 0 & 0,0000E+00 \\ \hline
complex1 & 14 & 4,7140E+05 & 16 & 4,0300E+00 & 15 & 4,9191E+05 & 37 & 1,0140E+01 \\ \hline
cr42 & 14 & 1,8040E+03 & 15 & 4,0000E-02 & 14 & 1,8040E+03 & 15 & 4,0000E-02 \\ \hline
crew & 0 & 8,5880E+03 & 12 & 1,6000E-01 & 0 & 8,8230E+03 & 12 & 1,6000E-01 \\ \hline
dano3mip & 12 & 1,0902E+06 & 34 & 1,3600E+01 & 12 & 3,4466E+06 & 32 & 7,8680E+01 \\ \hline
dbah00 & 0 & 2,8099E+05 & 43 & 1,2400E+00 & 0 & 4,8499E+06 & 32 & 1,0051E+02 \\ \hline
dbic1 & 0 & 1,8631E+06 & 45 & 2,8700E+01 & -1 & 0,0000E+00 & 0 & 0,0000E+00 \\ \hline
dbir1 & 0 & 2,5235E+06 & 24 & 6,8810E+01 & 0 & 5,7974E+06 & 23 & 2,8589E+02 \\ \hline
dbir2 & 0 & 2,8185E+06 & 26 & 8,4010E+01 & 0 & 6,7499E+06 & 26 & 3,9022E+02 \\ \hline
delfland & 0 & 3,4610E+03 & 46 & 4,0000E-02 & 0 & 1,2234E+04 & 46 & 8,0000E-02 \\ \hline
disp3,mps & 3 & 0,0000E+00 & 0 & 0,0000E+00 & 3 & 0,0000E+00 & 0 & 0,0000E+00 \\ \hline
ds90 & 0 & 9,9700E+04 & 28 & 3,7000E-01 & 0 & 1,1216E+06 & 21 & 9,9200E+00 \\ \hline
dsbmip & 0 & 1,6789E+04 & 48 & 1,6000E-01 & 0 & 5,9151E+04 & 48 & 4,3000E-01 \\ \hline
emsdz & 0 & 2,7497E+05 & 37 & 1,5000E+00 & 0 & 2,7832E+06 & 29 & 5,1400E+01 \\ \hline
f2177 & 0 & 1,2932E+07 & 2 & 2,8520E+01 & 0 & 2,5060E+07 & 2 & 1,0095E+02 \\ \hline
from-lp-file & -1 & 0,0000E+00 & 0 & 0,0000E+00 & -1 & 0,0000E+00 & 0 & 0,0000E+00 \\ \hline
gamsmod & 0 & 1,6610E+03 & 3 & 1,0000E-02 & 0 & 5,8360E+03 & 3 & 0,0000E+00 \\ \hline
iiasa & 0 & 4,5250E+03 & 15 & 3,0000E-02 & 0 & 4,6540E+03 & 15 & 4,0000E-02 \\ \hline
indata & 0 & 1,0008E+06 & 14 & 5,4600E+00 & 0 & 8,3833E+05 & 14 & 1,0150E+01 \\ \hline
jendrec1 & 6 & 0,0000E+00 & 0 & 0,0000E+00 & 6 & 0,0000E+00 & 0 & 0,0000E+00 \\ \hline
kl02,mps & 3 & 0,0000E+00 & 0 & 0,0000E+00 & 3 & 0,0000E+00 & 0 & 0,0000E+00 \\ \hline
kleemin3 & 0 & 6,0000E+00 & 5 & 0,0000E+00 & 0 & 6,0000E+00 & 5 & 0,0000E+00 \\ \hline
kleemin4 & 0 & 1,0000E+01 & 6 & 0,0000E+00 & 0 & 1,0000E+01 & 6 & 0,0000E+00 \\ \hline
kleemin5 & 0 & 1,5000E+01 & 8 & 0,0000E+00 & 0 & 1,5000E+01 & 8 & 0,0000E+00 \\ \hline
kleemin6 & 0 & 2,1000E+01 & 8 & 0,0000E+00 & 0 & 2,1000E+01 & 8 & 0,0000E+00 \\ \hline
kleemin7 & 0 & 2,8000E+01 & 8 & 0,0000E+00 & 0 & 2,8000E+01 & 8 & 0,0000E+00 \\ \hline
kleemin8 & 0 & 3,6000E+01 & 9 & 0,0000E+00 & 0 & 3,6000E+01 & 9 & 0,0000E+00 \\ \hline
l09a13l1d & 14 & 6,8450E+03 & 12 & 2,0000E-02 & 14 & 7,0910E+03 & 13 & 2,0000E-02 \\ \hline
leader & 0 & 3,9554E+05 & 49 & 2,8600E+00 & 14 & 4,7594E+06 & 48 & 5,1883E+02 \\ \hline
lindo & 0 & 2,0000E+01 & 9 & 0,0000E+00 & 0 & 2,0000E+01 & 9 & 0,0000E+00 \\ \hline
lp22 & 0 & 9,4273E+05 & 56 & 1,5690E+01 & 0 & 2,0149E+06 & 37 & 3,5400E+01 \\ \hline
lpl1 & 0 & 9,6124E+05 & 68 & 1,1760E+01 & 0 & 5,1342E+07 & 52 & 6,7786E+03 \\ \hline
lpl2 & 0 & 4,2869E+04 & 21 & 1,8000E-01 & 0 & 5,0591E+05 & 20 & 3,0700E+00 \\ \hline
lpl3 & 0 & 1,4890E+05 & 36 & 1,1200E+00 & 0 & 4,4995E+06 & 31 & 1,2124E+02 \\ \hline
model10 & 0 & 3,2144E+05 & 45 & 2,0800E+00 & 0 & 1,1938E+06 & 37 & 1,4640E+01 \\ \hline
model11 & 0 & 1,4007E+05 & 25 & 6,4000E-01 & 0 & 7,7321E+05 & 21 & 5,2500E+00 \\ \hline
model1 & 0 & 6,8610E+03 & 8 & 1,0000E-02 & 0 & 1,5093E+04 & 8 & 2,0000E-02 \\ \hline
model2 & 0 & 1,2167E+04 & 23 & 5,0000E-02 & 0 & 1,5749E+04 & 23 & 8,0000E-02 \\ \hline
model3 & 5 & 0,0000E+00 & 0 & 0,0000E+00 & 5 & 0,0000E+00 & 0 & 0,0000E+00 \\ \hline
model4 & 0 & 8,8743E+04 & 35 & 5,0000E-01 & 0 & 1,8571E+05 & 35 & 1,3200E+00 \\ \hline
model5 & 0 & 1,4052E+05 & 44 & 1,2300E+00 & 0 & 2,8807E+05 & 44 & 3,2300E+00 \\ \hline
model6 & 0 & 1,0634E+05 & 33 & 4,3000E-01 & 0 & 3,8909E+05 & 33 & 2,1400E+00 \\ \hline
model7 & 0 & 1,1000E+05 & 39 & 7,1000E-01 & 0 & 2,5601E+05 & 39 & 2,2000E+00 \\ \hline
model8 & 0 & 2,2376E+05 & 17 & 4,2000E-01 & 0 & 2,7580E+05 & 17 & 1,1400E+00 \\ \hline
model9 & 0 & 7,0839E+04 & 41 & 6,0000E-01 & 0 & 1,0504E+05 & 42 & 1,4400E+00 \\ \hline
nemsafm & 0 & 1,0160E+03 & 16 & 1,0000E-02 & 0 & 1,3240E+03 & 16 & 2,0000E-02 \\ \hline
nemscem & 0 & 2,5350E+03 & 18 & 2,0000E-02 & 0 & 1,4749E+04 & 18 & 3,0000E-02 \\ \hline
nemsemm1 & 0 & 2,4298E+05 & 57 & 1,5710E+01 & 0 & 3,9867E+05 & 57 & 1,9620E+01 \\ \hline
nemsemm2 & 0 & 6,2496E+04 & 35 & 8,7000E-01 & 0 & 3,8140E+05 & 35 & 3,3700E+00 \\ \hline
nemspmm1 & 0 & 1,1464E+05 & 39 & 7,4000E-01 & 0 & 8,6560E+05 & 31 & 9,0400E+00 \\ \hline
nemspmm2,in & 0 & 1,1890E+05 & 44 & 1,0100E+00 & 0 & 7,1490E+05 & 34 & 7,4900E+00 \\ \hline
nemswrld & 0 & 7,8312E+05 & 42 & 7,8000E+00 & 0 & 4,4720E+06 & 38 & 9,3030E+01 \\ \hline
nsct1 & 0 & 5,4180E+06 & 21 & 1,7080E+02 & 0 & 8,5935E+06 & 22 & 6,1142E+02 \\ \hline
nsct2 & 0 & 5,6456E+06 & 28 & 2,3305E+02 & 0 & 8,3502E+06 & 30 & 7,0599E+02 \\ \hline
nsic1 & 0 & 5,0420E+03 & 10 & 2,0000E-02 & 0 & 7,2400E+03 & 10 & 3,0000E-02 \\ \hline
nsic2 & 0 & 5,5160E+03 & 14 & 2,0000E-02 & 0 & 8,0110E+03 & 14 & 4,0000E-02 \\ \hline
nsir1 & 0 & 3,0346E+05 & 21 & 3,6600E+00 & 0 & 6,3039E+05 & 17 & 1,0390E+01 \\ \hline
nsir2 & 0 & 3,2795E+05 & 30 & 5,9400E+00 & 0 & 6,5743E+05 & 26 & 1,4420E+01 \\ \hline
nw14,mps & 3 & 0,0000E+00 & 0 & 0,0000E+00 & 3 & 0,0000E+00 & 0 & 0,0000E+00 \\ \hline
olivier & 14 & 4,6137E+05 & 50 & 4,5700E+00 & 14 & 5,6194E+06 & 43 & 4,4814E+02 \\ \hline
or4,mps & 3 & 0,0000E+00 & 0 & 0,0000E+00 & 3 & 0,0000E+00 & 0 & 0,0000E+00 \\ \hline
orna1,mps & 3 & 0,0000E+00 & 0 & 0,0000E+00 & 3 & 0,0000E+00 & 0 & 0,0000E+00 \\ \hline
orna2,mps & 3 & 0,0000E+00 & 0 & 0,0000E+00 & 3 & 0,0000E+00 & 0 & 0,0000E+00 \\ \hline
orna3,mps & 3 & 0,0000E+00 & 0 & 0,0000E+00 & 3 & 0,0000E+00 & 0 & 0,0000E+00 \\ \hline
orna7,mps & 3 & 0,0000E+00 & 0 & 0,0000E+00 & 3 & 0,0000E+00 & 0 & 0,0000E+00 \\ \hline
orswq2,mps & 3 & 0,0000E+00 & 0 & 0,0000E+00 & 3 & 0,0000E+00 & 0 & 0,0000E+00 \\ \hline
p0033 & 0 & 5,8000E+01 & 8 & 0,0000E+00 & 0 & 6,0000E+01 & 8 & 0,0000E+00 \\ \hline
p0040 & 0 & 9,6000E+01 & 9 & 0,0000E+00 & 0 & 9,6000E+01 & 9 & 0,0000E+00 \\ \hline
p0201 & 0 & 2,9220E+03 & 5 & 1,0000E-02 & 0 & 7,8160E+03 & 5 & 1,0000E-02 \\ \hline
p0282 & 0 & 9,5870E+03 & 13 & 2,0000E-02 & 0 & 1,3065E+04 & 13 & 4,0000E-02 \\ \hline
p0291 & 0 & 4,8890E+03 & 12 & 1,0000E-02 & 0 & 8,7050E+03 & 12 & 2,0000E-02 \\ \hline
p0548 & 0 & 1,0530E+03 & 24 & 2,0000E-02 & 0 & 1,1890E+03 & 24 & 2,0000E-02 \\ \hline
p05 & 0 & 2,3647E+05 & 23 & 7,1000E-01 & 0 & 2,2887E+06 & 19 & 2,3310E+01 \\ \hline
p10 & 0 & 4,6962E+05 & 27 & 1,6900E+00 & 0 & 7,7808E+06 & 29 & 2,5367E+02 \\ \hline
p12345 & 0 & 6,4360E+03 & 14 & 8,0000E-02 & 0 & 6,4360E+03 & 14 & 8,0000E-02 \\ \hline
p19328 & 0 & 1,6932E+04 & 15 & 2,7000E-01 & 0 & 1,6932E+04 & 15 & 3,1000E-01 \\ \hline
p19 & 0 & 2,9710E+03 & 19 & 2,0000E-02 & 0 & 9,8150E+03 & 19 & 3,0000E-02 \\ \hline
p2756 & 0 & 6,5240E+03 & 15 & 5,0000E-02 & 0 & 6,7656E+04 & 15 & 2,0000E-01 \\ \hline
pcb1000 & 0 & 2,9753E+04 & 23 & 1,8000E-01 & 0 & 3,9440E+05 & 23 & 1,8200E+00 \\ \hline
pcb3000 & 0 & 1,0109E+05 & 25 & 5,6000E-01 & 0 & 4,3551E+06 & 19 & 4,9320E+01 \\ \hline
primagaz & 0 & 1,2404E+04 & 13 & 1,1000E-01 & 0 & 1,2404E+04 & 13 & 1,4000E-01 \\ \hline
progas & 0 & 3,0946E+04 & 15 & 6,0000E-02 & 0 & 4,8144E+04 & 15 & 2,3000E-01 \\ \hline
ps & 15 & 1,1054E+07 & 40 & 5,1600E+02 & 15 & 1,2683E+07 & 48 & 6,2349E+02 \\ \hline
qiu & 0 & 4,6492E+04 & 5 & 2,0000E-02 & 0 & 1,0207E+05 & 5 & 8,0000E-02 \\ \hline
r05 & 0 & 4,3308E+05 & 23 & 1,7800E+00 & 0 & 2,1250E+06 & 18 & 2,6230E+01 \\ \hline
radio & -1 & 0,0000E+00 & 0 & 0,0000E+00 & -11 & 0,0000E+00 & 0 & 0,0000E+00 \\ \hline
rlf,pre,dual & 0 & 8,4078E+05 & 13 & 2,3400E+00 & 0 & 6,0951E+06 & 13 & 4,8841E+02 \\ \hline
routing & 0 & 3,1143E+06 & 20 & 6,0600E+00 & 0 & 1,2819E+07 & 17 & 8,7570E+01 \\ \hline
sc205 & 12 & 5,0422E+04 & 15 & 5,3000E-01 & 12 & 5,3625E+04 & 15 & 6,0000E-01 \\ \hline
seymour & 0 & 4,5141E+06 & 15 & 3,6330E+01 & 0 & 8,2165E+06 & 15 & 1,3423E+02 \\ \hline
slp-tsk & 0 & 4,0878E+06 & 19 & 6,4860E+01 & 0 & 4,0918E+06 & 19 & 1,2514E+02 \\ \hline
southern1 & 0 & 5,6710E+06 & 15 & 6,5000E+01 & 0 & 7,7658E+06 & 15 & 6,6628E+02 \\ \hline
sturing & 0 & 1,8861E+04 & 36 & 1,5000E-01 & 0 & 3,6899E+04 & 36 & 3,0000E-01 \\ \hline
sws & 0 & 5,8235E+04 & 10 & 2,1000E-01 & 0 & 4,6605E+05 & 10 & 5,9000E-01 \\ \hline
t0331-4l & 0 & 2,0167E+05 & 37 & 6,5800E+00 & 0 & 2,1556E+05 & 37 & 7,0600E+00 \\ \hline
test & 0 & 6,5540E+03 & 46 & 2,7000E-01 & 0 & 6,5540E+03 & 46 & 2,9000E-01 \\ \hline
testps,mod & 0 & 3,1000E+01 & 13 & 0,0000E+00 & 0 & 3,1000E+01 & 13 & 0,0000E+00 \\ \hline
unilever2 & 15 & 4,2276E+05 & 36 & 3,4700E+00 & 0 & 5,8050E+06 & 73 & 2,4664E+02 \\ \hline
us04,mps & 3 & 0,0000E+00 & 0 & 0,0000E+00 & 3 & 0,0000E+00 & 0 & 0,0000E+00 \\ \hline
world & 14 & 1,1112E+06 & 60 & 9,2800E+00 & -1 & 0,0000E+00 & 0 & 0,0000E+00 \\ \hline
zed & 0 & 3,3800E+03 & 9 & 2,0000E-02 & 0 & 3,9520E+03 & 9 & 2,0000E-02 \\ \hline
zz & 15 & 2,8570E+05 & 43 & 3,7000E+00 & 15 & 9,9260E+05 & 43 & 1,0700E+01 \\ \hline
\end{tabular}

\end{table}
\begin{table}
  \centering
  \caption{Resultados experimentais da biblioteca Netlib (1/3).}
  \label{tab:resulnet0}
  \begin{tabular}{|l|r|r|r|r|r|r|r|r|}
\hline
\multicolumn{1}{|c|}{Problema} & \multicolumn{4}{|c|}{MMD} &         \multicolumn{4}{|c|}{RCM} \\ \hline
Nome & R & NNZ & IT & T & R & NNZ & IT & T \\ \hline
25fv47 & 0 & 33809 & 25 & 1,100000e-01 & 0 & 70740 & 25 & 2,800000e-01 \\ \hline
80bau3b & 0 & 41367 & 36 & 3,100000e-01 & 0 & 230301 & 36 & 1,440000e+00 \\ \hline
adlittle & 0 & 404 & 11 & 0,000000e+00 & 0 & 503 & 11 & 0,000000e+00 \\ \hline
afiro & 0 & 107 & 7 & 0,000000e+00 & 0 & 166 & 7 & 0,000000e+00 \\ \hline
agg2 & 0 & 21482 & 21 & 6,000000e-02 & 0 & 59242 & 21 & 1,700000e-01 \\ \hline
agg3 & 0 & 21482 & 19 & 6,000000e-02 & 0 & 59242 & 19 & 1,500000e-01 \\ \hline
agg & 0 & 12297 & 18 & 3,000000e-02 & 0 & 26647 & 18 & 6,000000e-02 \\ \hline
bandm & 0 & 3936 & 16 & 2,000000e-02 & 0 & 7083 & 16 & 2,000000e-02 \\ \hline
beaconfd & 0 & 820 & 10 & 1,000000e-02 & 0 & 1029 & 10 & 0,000000e+00 \\ \hline
blend & 0 & 913 & 9 & 0,000000e+00 & 0 & 1813 & 9 & 1,000000e-02 \\ \hline
bnl1 & 0 & 12089 & 33 & 6,000000e-02 & 0 & 31463 & 33 & 1,400000e-01 \\ \hline
bnl2 & 0 & 81275 & 33 & 3,100000e-01 & 0 & 233939 & 33 & 1,340000e+00 \\ \hline
boeing1 & 0 & 5725 & 18 & 2,000000e-02 & 0 & 39025 & 19 & 9,000000e-02 \\ \hline
boeing2 & 0 & 2029 & 13 & 1,000000e-02 & 0 & 3093 & 13 & 2,000000e-02 \\ \hline
bore3d & 0 & 1034 & 15 & 1,000000e-02 & 0 & 1408 & 15 & 1,000000e-02 \\ \hline
brandy & 14 & 2755 & 17 & 2,000000e-02 & 14 & 3708 & 17 & 2,000000e-02 \\ \hline
capri & 0 & 3962 & 17 & 2,000000e-02 & 0 & 5711 & 17 & 2,000000e-02 \\ \hline
cycle & 0 & 56102 & 23 & 1,600000e-01 & 0 & 178165 & 23 & 8,000000e-01 \\ \hline
czprob & 0 & 3520 & 26 & 4,000000e-02 & 0 & 91320 & 26 & 2,200000e-01 \\ \hline
d2q06c & 0 & 137349 & 27 & 5,400000e-01 & 0 & 455430 & 24 & 3,010000e+00 \\ \hline
d6cube & 0 & 54840 & 17 & 2,200000e-01 & 0 & 65467 & 17 & 2,800000e-01 \\ \hline
degen2 & 0 & 16319 & 11 & 3,000000e-02 & 0 & 40896 & 11 & 8,000000e-02 \\ \hline
degen3 & 0 & 120906 & 16 & 5,000000e-01 & 0 & 531478 & 13 & 2,530000e+00 \\ \hline
dfl001 & 0 & 1638085 & 45 & 2,443000e+01 & 15 & 5720111 & 60 & 3,470800e+02 \\ \hline
e226 & 0 & 3229 & 18 & 2,000000e-02 & 0 & 7715 & 18 & 3,000000e-02 \\ \hline
etamacro & 0 & 10843 & 26 & 3,000000e-02 & 0 & 23260 & 26 & 7,000000e-02 \\ \hline
fffff800 & 0 & 9573 & 29 & 7,000000e-02 & 0 & 33386 & 29 & 1,300000e-01 \\ \hline
finnis & 0 & 4984 & 22 & 3,000000e-02 & 0 & 10466 & 22 & 4,000000e-02 \\ \hline
fit1d & 0 & 296 & 17 & 6,000000e-02 & 0 & 298 & 17 & 6,000000e-02 \\ \hline
fit1p & 0 & 627 & 16 & 7,000000e-02 & 0 & 627 & 16 & 8,000000e-02 \\ \hline
fit2d & 0 & 324 & 22 & 7,100000e-01 & 0 & 325 & 22 & 7,100000e-01 \\ \hline
fit2p & 0 & 3000 & 20 & 5,400000e-01 & 0 & 3000 & 20 & 5,800000e-01 \\ \hline
forplan & 0 & 3304 & 21 & 3,000000e-02 & 0 & 4352 & 21 & 3,000000e-02 \\ \hline
ganges & 0 & 29677 & 16 & 6,000000e-02 & 0 & 48834 & 16 & 1,600000e-01 \\ \hline
gfrd-pnc & 0 & 2112 & 16 & 2,000000e-02 & 0 & 3316 & 16 & 2,000000e-02 \\ \hline
greenbea & 14 & 49055 & 50 & 4,400000e-01 & 14 & 150868 & 49 & 1,520000e+00 \\ \hline
greenbeb & 14 & 47783 & 40 & 3,500000e-01 & 14 & 167691 & 39 & 1,180000e+00 \\ \hline
grow15 & 0 & 6090 & 19 & 4,000000e-02 & 0 & 6090 & 19 & 5,000000e-02 \\ \hline
grow22 & 0 & 9058 & 21 & 6,000000e-02 & 0 & 9030 & 21 & 8,000000e-02 \\ \hline
grow7 & 0 & 2730 & 16 & 2,000000e-02 & 0 & 2730 & 16 & 2,000000e-02 \\ \hline
israel & 0 & 11488 & 19 & 4,000000e-02 & 0 & 11506 & 19 & 6,000000e-02 \\ \hline
kb2 & 0 & 503 & 11 & 0,000000e+00 & 0 & 747 & 11 & 0,000000e+00 \\ \hline
lotfi & 0 & 1369 & 13 & 1,000000e-02 & 0 & 2694 & 13 & 1,000000e-02 \\ \hline
maros & 0 & 13454 & 18 & 4,000000e-02 & 0 & 33456 & 18 & 1,000000e-01 \\ \hline
maros-r7 & 0 & 534188 & 14 & 1,610000e+00 & 0 & 439842 & 14 & 4,100000e+00 \\ \hline
modszk1 & 0 & 10550 & 20 & 3,000000e-02 & 0 & 82365 & 20 & 1,700000e-01 \\ \hline
nesm & 0 & 21776 & 25 & 1,200000e-01 & 0 & 29048 & 25 & 1,600000e-01 \\ \hline
perold & 0 & 21782 & 33 & 9,000000e-02 & 0 & 44505 & 33 & 2,100000e-01 \\ \hline
pilot4 & 0 & 14279 & 46 & 1,400000e-01 & 0 & 26482 & 46 & 2,100000e-01 \\ \hline
pilot87 & 0 & 425654 & 29 & 3,350000e+00 & 0 & 771344 & 25 & 7,710000e+00 \\ \hline
pilot,ja & 0 & 47924 & 29 & 2,000000e-01 & 0 & 95733 & 29 & 4,800000e-01 \\ \hline
pilot & 0 & 200812 & 34 & 1,300000e+00 & 0 & 435889 & 30 & 4,180000e+00 \\ \hline
pilotnov & 0 & 46353 & 16 & 1,100000e-01 & 0 & 95971 & 16 & 2,600000e-01 \\ \hline
pilot,we & 0 & 15605 & 45 & 1,300000e-01 & 0 & 32919 & 45 & 2,100000e-01 \\ \hline
recipe & 0 & 277 & 8 & 0,000000e+00 & 0 & 278 & 8 & 1,000000e-02 \\ \hline
sc105 & 0 & 569 & 9 & 0,000000e+00 & 0 & 825 & 9 & 0,000000e+00 \\ \hline
sc205 & 0 & 1156 & 10 & 1,000000e-02 & 0 & 1716 & 10 & 1,000000e-02 \\ \hline
sc50a & 0 & 242 & 7 & 0,000000e+00 & 0 & 326 & 7 & 0,000000e+00 \\ \hline
sc50b & 0 & 231 & 6 & 0,000000e+00 & 0 & 361 & 6 & 0,000000e+00 \\ \hline
scagr25 & 0 & 2944 & 17 & 2,000000e-02 & 0 & 4258 & 17 & 3,000000e-02 \\ \hline
scagr7 & 0 & 730 & 13 & 1,000000e-02 & 0 & 1036 & 13 & 1,000000e-02 \\ \hline
scfxm1 & 0 & 4430 & 17 & 2,000000e-02 & 0 & 7419 & 17 & 3,000000e-02 \\ \hline
scfxm2 & 14 & 9284 & 19 & 4,000000e-02 & 14 & 15458 & 19 & 7,000000e-02 \\ \hline
scfxm3 & 0 & 14138 & 19 & 6,000000e-02 & 0 & 23442 & 19 & 1,300000e-01 \\ \hline
scorpion & 0 & 2275 & 11 & 1,000000e-02 & 0 & 4926 & 11 & 2,000000e-02 \\ \hline
scrs8 & 0 & 4477 & 21 & 2,000000e-02 & 0 & 8304 & 21 & 3,000000e-02 \\ \hline
scsd1 & 0 & 1392 & 8 & 1,000000e-02 & 0 & 1472 & 8 & 0,000000e+00 \\ \hline
scsd6 & 0 & 2545 & 11 & 2,000000e-02 & 0 & 2740 & 11 & 2,000000e-02 \\ \hline
scsd8 & 0 & 5879 & 10 & 3,000000e-02 & 0 & 5942 & 10 & 3,000000e-02 \\ \hline
sctap1 & 0 & 2435 & 14 & 1,000000e-02 & 0 & 5367 & 14 & 2,000000e-02 \\ \hline
sctap2 & 0 & 11736 & 12 & 3,000000e-02 & 0 & 55394 & 12 & 1,100000e-01 \\ \hline
sctap3 & 0 & 16100 & 14 & 6,000000e-02 & 0 & 102574 & 14 & 2,600000e-01 \\ \hline
seba & 0 & 53711 & 13 & 1,700000e-01 & 0 & 55513 & 13 & 4,000000e-01 \\ \hline
share1b & 0 & 1417 & 18 & 1,000000e-02 & 0 & 1762 & 18 & 2,000000e-02 \\ \hline
share2b & 0 & 1041 & 17 & 1,000000e-02 & 0 & 1542 & 17 & 1,000000e-02 \\ \hline
shell & 0 & 3983 & 20 & 2,000000e-02 & 0 & 9539 & 20 & 4,000000e-02 \\ \hline
ship04l & 0 & 2641 & 12 & 2,000000e-02 & 0 & 7944 & 12 & 2,000000e-02 \\ \hline
ship04s & 0 & 1778 & 12 & 1,000000e-02 & 0 & 4091 & 12 & 2,000000e-02 \\ \hline
ship08l & 0 & 4442 & 14 & 3,000000e-02 & 0 & 15947 & 14 & 4,000000e-02 \\ \hline
ship08s & 0 & 2295 & 11 & 2,000000e-02 & 0 & 4290 & 11 & 2,000000e-02 \\ \hline
ship12l & 0 & 5506 & 15 & 4,000000e-02 & 0 & 16980 & 15 & 5,000000e-02 \\ \hline
ship12s & 0 & 2507 & 12 & 2,000000e-02 & 0 & 3801 & 12 & 2,000000e-02 \\ \hline
sierra & 0 & 12862 & 18 & 6,000000e-02 & 0 & 43552 & 18 & 1,400000e-01 \\ \hline
stair & 0 & 14682 & 13 & 3,000000e-02 & 0 & 15047 & 13 & 5,000000e-02 \\ \hline
standata & 0 & 2395 & 13 & 1,000000e-02 & 0 & 3953 & 13 & 2,000000e-02 \\ \hline
standgub & 0 & 2395 & 13 & 1,000000e-02 & 0 & 3953 & 13 & 1,000000e-02 \\ \hline
standmps & 0 & 3957 & 23 & 2,000000e-02 & 0 & 9136 & 23 & 3,000000e-02 \\ \hline
stocfor1 & 0 & 805 & 11 & 0,000000e+00 & 0 & 1388 & 11 & 1,000000e-02 \\ \hline
stocfor2 & 0 & 22841 & 19 & 7,000000e-02 & 0 & 32610 & 19 & 2,400000e-01 \\ \hline
tuff & 0 & 7051 & 18 & 2,000000e-02 & 0 & 9820 & 18 & 4,000000e-02 \\ \hline
vtp,base & 0 & 505 & 10 & 0,000000e+00 & 0 & 504 & 10 & 0,000000e+00 \\ \hline
wood1p & 0 & 11645 & 22 & 3,400000e-01 & 0 & 14537 & 22 & 3,600000e-01 \\ \hline
woodw & 0 & 30027 & 30 & 1,800000e-01 & 0 & 126724 & 30 & 5,800000e-01 \\ \hline
\end{tabular}

\end{table}
\begin{table}
  \centering
  \caption{Resultados experimentais da biblioteca Netlib (2/3).}
  \label{tab:resulnet1}
  \begin{tabular}{|l|r|r|r|r|r|r|r|r|}
\hline
\multicolumn{1}{|c|}{Problema} & \multicolumn{4}{|c|}{MMD} &         \multicolumn{4}{|c|}{RCM} \\ \hline
Nome & R & NNZ & IT & T & R & NNZ & IT & T \\ \hline
25fv47 & 0 & 33809 & 25 & 1,100000e-01 & 0 & 70740 & 25 & 2,800000e-01 \\ \hline
80bau3b & 0 & 41367 & 36 & 3,100000e-01 & 0 & 230301 & 36 & 1,440000e+00 \\ \hline
adlittle & 0 & 404 & 11 & 0,000000e+00 & 0 & 503 & 11 & 0,000000e+00 \\ \hline
afiro & 0 & 107 & 7 & 0,000000e+00 & 0 & 166 & 7 & 0,000000e+00 \\ \hline
agg2 & 0 & 21482 & 21 & 6,000000e-02 & 0 & 59242 & 21 & 1,700000e-01 \\ \hline
agg3 & 0 & 21482 & 19 & 6,000000e-02 & 0 & 59242 & 19 & 1,500000e-01 \\ \hline
agg & 0 & 12297 & 18 & 3,000000e-02 & 0 & 26647 & 18 & 6,000000e-02 \\ \hline
bandm & 0 & 3936 & 16 & 2,000000e-02 & 0 & 7083 & 16 & 2,000000e-02 \\ \hline
beaconfd & 0 & 820 & 10 & 1,000000e-02 & 0 & 1029 & 10 & 0,000000e+00 \\ \hline
blend & 0 & 913 & 9 & 0,000000e+00 & 0 & 1813 & 9 & 1,000000e-02 \\ \hline
bnl1 & 0 & 12089 & 33 & 6,000000e-02 & 0 & 31463 & 33 & 1,400000e-01 \\ \hline
bnl2 & 0 & 81275 & 33 & 3,100000e-01 & 0 & 233939 & 33 & 1,340000e+00 \\ \hline
boeing1 & 0 & 5725 & 18 & 2,000000e-02 & 0 & 39025 & 19 & 9,000000e-02 \\ \hline
boeing2 & 0 & 2029 & 13 & 1,000000e-02 & 0 & 3093 & 13 & 2,000000e-02 \\ \hline
bore3d & 0 & 1034 & 15 & 1,000000e-02 & 0 & 1408 & 15 & 1,000000e-02 \\ \hline
brandy & 14 & 2755 & 17 & 2,000000e-02 & 14 & 3708 & 17 & 2,000000e-02 \\ \hline
capri & 0 & 3962 & 17 & 2,000000e-02 & 0 & 5711 & 17 & 2,000000e-02 \\ \hline
cycle & 0 & 56102 & 23 & 1,600000e-01 & 0 & 178165 & 23 & 8,000000e-01 \\ \hline
czprob & 0 & 3520 & 26 & 4,000000e-02 & 0 & 91320 & 26 & 2,200000e-01 \\ \hline
d2q06c & 0 & 137349 & 27 & 5,400000e-01 & 0 & 455430 & 24 & 3,010000e+00 \\ \hline
d6cube & 0 & 54840 & 17 & 2,200000e-01 & 0 & 65467 & 17 & 2,800000e-01 \\ \hline
degen2 & 0 & 16319 & 11 & 3,000000e-02 & 0 & 40896 & 11 & 8,000000e-02 \\ \hline
degen3 & 0 & 120906 & 16 & 5,000000e-01 & 0 & 531478 & 13 & 2,530000e+00 \\ \hline
dfl001 & 0 & 1638085 & 45 & 2,443000e+01 & 15 & 5720111 & 60 & 3,470800e+02 \\ \hline
e226 & 0 & 3229 & 18 & 2,000000e-02 & 0 & 7715 & 18 & 3,000000e-02 \\ \hline
etamacro & 0 & 10843 & 26 & 3,000000e-02 & 0 & 23260 & 26 & 7,000000e-02 \\ \hline
fffff800 & 0 & 9573 & 29 & 7,000000e-02 & 0 & 33386 & 29 & 1,300000e-01 \\ \hline
finnis & 0 & 4984 & 22 & 3,000000e-02 & 0 & 10466 & 22 & 4,000000e-02 \\ \hline
fit1d & 0 & 296 & 17 & 6,000000e-02 & 0 & 298 & 17 & 6,000000e-02 \\ \hline
fit1p & 0 & 627 & 16 & 7,000000e-02 & 0 & 627 & 16 & 8,000000e-02 \\ \hline
fit2d & 0 & 324 & 22 & 7,100000e-01 & 0 & 325 & 22 & 7,100000e-01 \\ \hline
fit2p & 0 & 3000 & 20 & 5,400000e-01 & 0 & 3000 & 20 & 5,800000e-01 \\ \hline
forplan & 0 & 3304 & 21 & 3,000000e-02 & 0 & 4352 & 21 & 3,000000e-02 \\ \hline
ganges & 0 & 29677 & 16 & 6,000000e-02 & 0 & 48834 & 16 & 1,600000e-01 \\ \hline
gfrd-pnc & 0 & 2112 & 16 & 2,000000e-02 & 0 & 3316 & 16 & 2,000000e-02 \\ \hline
greenbea & 14 & 49055 & 50 & 4,400000e-01 & 14 & 150868 & 49 & 1,520000e+00 \\ \hline
greenbeb & 14 & 47783 & 40 & 3,500000e-01 & 14 & 167691 & 39 & 1,180000e+00 \\ \hline
grow15 & 0 & 6090 & 19 & 4,000000e-02 & 0 & 6090 & 19 & 5,000000e-02 \\ \hline
grow22 & 0 & 9058 & 21 & 6,000000e-02 & 0 & 9030 & 21 & 8,000000e-02 \\ \hline
grow7 & 0 & 2730 & 16 & 2,000000e-02 & 0 & 2730 & 16 & 2,000000e-02 \\ \hline
israel & 0 & 11488 & 19 & 4,000000e-02 & 0 & 11506 & 19 & 6,000000e-02 \\ \hline
kb2 & 0 & 503 & 11 & 0,000000e+00 & 0 & 747 & 11 & 0,000000e+00 \\ \hline
lotfi & 0 & 1369 & 13 & 1,000000e-02 & 0 & 2694 & 13 & 1,000000e-02 \\ \hline
maros & 0 & 13454 & 18 & 4,000000e-02 & 0 & 33456 & 18 & 1,000000e-01 \\ \hline
maros-r7 & 0 & 534188 & 14 & 1,610000e+00 & 0 & 439842 & 14 & 4,100000e+00 \\ \hline
modszk1 & 0 & 10550 & 20 & 3,000000e-02 & 0 & 82365 & 20 & 1,700000e-01 \\ \hline
nesm & 0 & 21776 & 25 & 1,200000e-01 & 0 & 29048 & 25 & 1,600000e-01 \\ \hline
perold & 0 & 21782 & 33 & 9,000000e-02 & 0 & 44505 & 33 & 2,100000e-01 \\ \hline
pilot4 & 0 & 14279 & 46 & 1,400000e-01 & 0 & 26482 & 46 & 2,100000e-01 \\ \hline
pilot87 & 0 & 425654 & 29 & 3,350000e+00 & 0 & 771344 & 25 & 7,710000e+00 \\ \hline
pilot,ja & 0 & 47924 & 29 & 2,000000e-01 & 0 & 95733 & 29 & 4,800000e-01 \\ \hline
pilot & 0 & 200812 & 34 & 1,300000e+00 & 0 & 435889 & 30 & 4,180000e+00 \\ \hline
pilotnov & 0 & 46353 & 16 & 1,100000e-01 & 0 & 95971 & 16 & 2,600000e-01 \\ \hline
pilot,we & 0 & 15605 & 45 & 1,300000e-01 & 0 & 32919 & 45 & 2,100000e-01 \\ \hline
recipe & 0 & 277 & 8 & 0,000000e+00 & 0 & 278 & 8 & 1,000000e-02 \\ \hline
sc105 & 0 & 569 & 9 & 0,000000e+00 & 0 & 825 & 9 & 0,000000e+00 \\ \hline
sc205 & 0 & 1156 & 10 & 1,000000e-02 & 0 & 1716 & 10 & 1,000000e-02 \\ \hline
sc50a & 0 & 242 & 7 & 0,000000e+00 & 0 & 326 & 7 & 0,000000e+00 \\ \hline
sc50b & 0 & 231 & 6 & 0,000000e+00 & 0 & 361 & 6 & 0,000000e+00 \\ \hline
scagr25 & 0 & 2944 & 17 & 2,000000e-02 & 0 & 4258 & 17 & 3,000000e-02 \\ \hline
scagr7 & 0 & 730 & 13 & 1,000000e-02 & 0 & 1036 & 13 & 1,000000e-02 \\ \hline
scfxm1 & 0 & 4430 & 17 & 2,000000e-02 & 0 & 7419 & 17 & 3,000000e-02 \\ \hline
scfxm2 & 14 & 9284 & 19 & 4,000000e-02 & 14 & 15458 & 19 & 7,000000e-02 \\ \hline
scfxm3 & 0 & 14138 & 19 & 6,000000e-02 & 0 & 23442 & 19 & 1,300000e-01 \\ \hline
scorpion & 0 & 2275 & 11 & 1,000000e-02 & 0 & 4926 & 11 & 2,000000e-02 \\ \hline
scrs8 & 0 & 4477 & 21 & 2,000000e-02 & 0 & 8304 & 21 & 3,000000e-02 \\ \hline
scsd1 & 0 & 1392 & 8 & 1,000000e-02 & 0 & 1472 & 8 & 0,000000e+00 \\ \hline
scsd6 & 0 & 2545 & 11 & 2,000000e-02 & 0 & 2740 & 11 & 2,000000e-02 \\ \hline
scsd8 & 0 & 5879 & 10 & 3,000000e-02 & 0 & 5942 & 10 & 3,000000e-02 \\ \hline
sctap1 & 0 & 2435 & 14 & 1,000000e-02 & 0 & 5367 & 14 & 2,000000e-02 \\ \hline
sctap2 & 0 & 11736 & 12 & 3,000000e-02 & 0 & 55394 & 12 & 1,100000e-01 \\ \hline
sctap3 & 0 & 16100 & 14 & 6,000000e-02 & 0 & 102574 & 14 & 2,600000e-01 \\ \hline
seba & 0 & 53711 & 13 & 1,700000e-01 & 0 & 55513 & 13 & 4,000000e-01 \\ \hline
share1b & 0 & 1417 & 18 & 1,000000e-02 & 0 & 1762 & 18 & 2,000000e-02 \\ \hline
share2b & 0 & 1041 & 17 & 1,000000e-02 & 0 & 1542 & 17 & 1,000000e-02 \\ \hline
shell & 0 & 3983 & 20 & 2,000000e-02 & 0 & 9539 & 20 & 4,000000e-02 \\ \hline
ship04l & 0 & 2641 & 12 & 2,000000e-02 & 0 & 7944 & 12 & 2,000000e-02 \\ \hline
ship04s & 0 & 1778 & 12 & 1,000000e-02 & 0 & 4091 & 12 & 2,000000e-02 \\ \hline
ship08l & 0 & 4442 & 14 & 3,000000e-02 & 0 & 15947 & 14 & 4,000000e-02 \\ \hline
ship08s & 0 & 2295 & 11 & 2,000000e-02 & 0 & 4290 & 11 & 2,000000e-02 \\ \hline
ship12l & 0 & 5506 & 15 & 4,000000e-02 & 0 & 16980 & 15 & 5,000000e-02 \\ \hline
ship12s & 0 & 2507 & 12 & 2,000000e-02 & 0 & 3801 & 12 & 2,000000e-02 \\ \hline
sierra & 0 & 12862 & 18 & 6,000000e-02 & 0 & 43552 & 18 & 1,400000e-01 \\ \hline
stair & 0 & 14682 & 13 & 3,000000e-02 & 0 & 15047 & 13 & 5,000000e-02 \\ \hline
standata & 0 & 2395 & 13 & 1,000000e-02 & 0 & 3953 & 13 & 2,000000e-02 \\ \hline
standgub & 0 & 2395 & 13 & 1,000000e-02 & 0 & 3953 & 13 & 1,000000e-02 \\ \hline
standmps & 0 & 3957 & 23 & 2,000000e-02 & 0 & 9136 & 23 & 3,000000e-02 \\ \hline
stocfor1 & 0 & 805 & 11 & 0,000000e+00 & 0 & 1388 & 11 & 1,000000e-02 \\ \hline
stocfor2 & 0 & 22841 & 19 & 7,000000e-02 & 0 & 32610 & 19 & 2,400000e-01 \\ \hline
tuff & 0 & 7051 & 18 & 2,000000e-02 & 0 & 9820 & 18 & 4,000000e-02 \\ \hline
vtp,base & 0 & 505 & 10 & 0,000000e+00 & 0 & 504 & 10 & 0,000000e+00 \\ \hline
wood1p & 0 & 11645 & 22 & 3,400000e-01 & 0 & 14537 & 22 & 3,600000e-01 \\ \hline
woodw & 0 & 30027 & 30 & 1,800000e-01 & 0 & 126724 & 30 & 5,800000e-01 \\ \hline
\end{tabular}

\end{table}
\begin{table}
  \centering
  \caption{Resultados experimentais da biblioteca Netlib (3/3).}
  \label{tab:resulnet2}
  \begin{tabular}{|l|r|r|r|r|r|r|r|r|}
\hline
\multicolumn{1}{|c|}{Problema} & \multicolumn{4}{|c|}{MMD} &         \multicolumn{4}{|c|}{RCM} \\ \hline
Nome & R & NNZ & IT & T & R & NNZ & IT & T \\ \hline
25fv47 & 0 & 33809 & 25 & 1,100000e-01 & 0 & 70740 & 25 & 2,800000e-01 \\ \hline
80bau3b & 0 & 41367 & 36 & 3,100000e-01 & 0 & 230301 & 36 & 1,440000e+00 \\ \hline
adlittle & 0 & 404 & 11 & 0,000000e+00 & 0 & 503 & 11 & 0,000000e+00 \\ \hline
afiro & 0 & 107 & 7 & 0,000000e+00 & 0 & 166 & 7 & 0,000000e+00 \\ \hline
agg2 & 0 & 21482 & 21 & 6,000000e-02 & 0 & 59242 & 21 & 1,700000e-01 \\ \hline
agg3 & 0 & 21482 & 19 & 6,000000e-02 & 0 & 59242 & 19 & 1,500000e-01 \\ \hline
agg & 0 & 12297 & 18 & 3,000000e-02 & 0 & 26647 & 18 & 6,000000e-02 \\ \hline
bandm & 0 & 3936 & 16 & 2,000000e-02 & 0 & 7083 & 16 & 2,000000e-02 \\ \hline
beaconfd & 0 & 820 & 10 & 1,000000e-02 & 0 & 1029 & 10 & 0,000000e+00 \\ \hline
blend & 0 & 913 & 9 & 0,000000e+00 & 0 & 1813 & 9 & 1,000000e-02 \\ \hline
bnl1 & 0 & 12089 & 33 & 6,000000e-02 & 0 & 31463 & 33 & 1,400000e-01 \\ \hline
bnl2 & 0 & 81275 & 33 & 3,100000e-01 & 0 & 233939 & 33 & 1,340000e+00 \\ \hline
boeing1 & 0 & 5725 & 18 & 2,000000e-02 & 0 & 39025 & 19 & 9,000000e-02 \\ \hline
boeing2 & 0 & 2029 & 13 & 1,000000e-02 & 0 & 3093 & 13 & 2,000000e-02 \\ \hline
bore3d & 0 & 1034 & 15 & 1,000000e-02 & 0 & 1408 & 15 & 1,000000e-02 \\ \hline
brandy & 14 & 2755 & 17 & 2,000000e-02 & 14 & 3708 & 17 & 2,000000e-02 \\ \hline
capri & 0 & 3962 & 17 & 2,000000e-02 & 0 & 5711 & 17 & 2,000000e-02 \\ \hline
cycle & 0 & 56102 & 23 & 1,600000e-01 & 0 & 178165 & 23 & 8,000000e-01 \\ \hline
czprob & 0 & 3520 & 26 & 4,000000e-02 & 0 & 91320 & 26 & 2,200000e-01 \\ \hline
d2q06c & 0 & 137349 & 27 & 5,400000e-01 & 0 & 455430 & 24 & 3,010000e+00 \\ \hline
d6cube & 0 & 54840 & 17 & 2,200000e-01 & 0 & 65467 & 17 & 2,800000e-01 \\ \hline
degen2 & 0 & 16319 & 11 & 3,000000e-02 & 0 & 40896 & 11 & 8,000000e-02 \\ \hline
degen3 & 0 & 120906 & 16 & 5,000000e-01 & 0 & 531478 & 13 & 2,530000e+00 \\ \hline
dfl001 & 0 & 1638085 & 45 & 2,443000e+01 & 15 & 5720111 & 60 & 3,470800e+02 \\ \hline
e226 & 0 & 3229 & 18 & 2,000000e-02 & 0 & 7715 & 18 & 3,000000e-02 \\ \hline
etamacro & 0 & 10843 & 26 & 3,000000e-02 & 0 & 23260 & 26 & 7,000000e-02 \\ \hline
fffff800 & 0 & 9573 & 29 & 7,000000e-02 & 0 & 33386 & 29 & 1,300000e-01 \\ \hline
finnis & 0 & 4984 & 22 & 3,000000e-02 & 0 & 10466 & 22 & 4,000000e-02 \\ \hline
fit1d & 0 & 296 & 17 & 6,000000e-02 & 0 & 298 & 17 & 6,000000e-02 \\ \hline
fit1p & 0 & 627 & 16 & 7,000000e-02 & 0 & 627 & 16 & 8,000000e-02 \\ \hline
fit2d & 0 & 324 & 22 & 7,100000e-01 & 0 & 325 & 22 & 7,100000e-01 \\ \hline
fit2p & 0 & 3000 & 20 & 5,400000e-01 & 0 & 3000 & 20 & 5,800000e-01 \\ \hline
forplan & 0 & 3304 & 21 & 3,000000e-02 & 0 & 4352 & 21 & 3,000000e-02 \\ \hline
ganges & 0 & 29677 & 16 & 6,000000e-02 & 0 & 48834 & 16 & 1,600000e-01 \\ \hline
gfrd-pnc & 0 & 2112 & 16 & 2,000000e-02 & 0 & 3316 & 16 & 2,000000e-02 \\ \hline
greenbea & 14 & 49055 & 50 & 4,400000e-01 & 14 & 150868 & 49 & 1,520000e+00 \\ \hline
greenbeb & 14 & 47783 & 40 & 3,500000e-01 & 14 & 167691 & 39 & 1,180000e+00 \\ \hline
grow15 & 0 & 6090 & 19 & 4,000000e-02 & 0 & 6090 & 19 & 5,000000e-02 \\ \hline
grow22 & 0 & 9058 & 21 & 6,000000e-02 & 0 & 9030 & 21 & 8,000000e-02 \\ \hline
grow7 & 0 & 2730 & 16 & 2,000000e-02 & 0 & 2730 & 16 & 2,000000e-02 \\ \hline
israel & 0 & 11488 & 19 & 4,000000e-02 & 0 & 11506 & 19 & 6,000000e-02 \\ \hline
kb2 & 0 & 503 & 11 & 0,000000e+00 & 0 & 747 & 11 & 0,000000e+00 \\ \hline
lotfi & 0 & 1369 & 13 & 1,000000e-02 & 0 & 2694 & 13 & 1,000000e-02 \\ \hline
maros & 0 & 13454 & 18 & 4,000000e-02 & 0 & 33456 & 18 & 1,000000e-01 \\ \hline
maros-r7 & 0 & 534188 & 14 & 1,610000e+00 & 0 & 439842 & 14 & 4,100000e+00 \\ \hline
modszk1 & 0 & 10550 & 20 & 3,000000e-02 & 0 & 82365 & 20 & 1,700000e-01 \\ \hline
nesm & 0 & 21776 & 25 & 1,200000e-01 & 0 & 29048 & 25 & 1,600000e-01 \\ \hline
perold & 0 & 21782 & 33 & 9,000000e-02 & 0 & 44505 & 33 & 2,100000e-01 \\ \hline
pilot4 & 0 & 14279 & 46 & 1,400000e-01 & 0 & 26482 & 46 & 2,100000e-01 \\ \hline
pilot87 & 0 & 425654 & 29 & 3,350000e+00 & 0 & 771344 & 25 & 7,710000e+00 \\ \hline
pilot,ja & 0 & 47924 & 29 & 2,000000e-01 & 0 & 95733 & 29 & 4,800000e-01 \\ \hline
pilot & 0 & 200812 & 34 & 1,300000e+00 & 0 & 435889 & 30 & 4,180000e+00 \\ \hline
pilotnov & 0 & 46353 & 16 & 1,100000e-01 & 0 & 95971 & 16 & 2,600000e-01 \\ \hline
pilot,we & 0 & 15605 & 45 & 1,300000e-01 & 0 & 32919 & 45 & 2,100000e-01 \\ \hline
recipe & 0 & 277 & 8 & 0,000000e+00 & 0 & 278 & 8 & 1,000000e-02 \\ \hline
sc105 & 0 & 569 & 9 & 0,000000e+00 & 0 & 825 & 9 & 0,000000e+00 \\ \hline
sc205 & 0 & 1156 & 10 & 1,000000e-02 & 0 & 1716 & 10 & 1,000000e-02 \\ \hline
sc50a & 0 & 242 & 7 & 0,000000e+00 & 0 & 326 & 7 & 0,000000e+00 \\ \hline
sc50b & 0 & 231 & 6 & 0,000000e+00 & 0 & 361 & 6 & 0,000000e+00 \\ \hline
scagr25 & 0 & 2944 & 17 & 2,000000e-02 & 0 & 4258 & 17 & 3,000000e-02 \\ \hline
scagr7 & 0 & 730 & 13 & 1,000000e-02 & 0 & 1036 & 13 & 1,000000e-02 \\ \hline
scfxm1 & 0 & 4430 & 17 & 2,000000e-02 & 0 & 7419 & 17 & 3,000000e-02 \\ \hline
scfxm2 & 14 & 9284 & 19 & 4,000000e-02 & 14 & 15458 & 19 & 7,000000e-02 \\ \hline
scfxm3 & 0 & 14138 & 19 & 6,000000e-02 & 0 & 23442 & 19 & 1,300000e-01 \\ \hline
scorpion & 0 & 2275 & 11 & 1,000000e-02 & 0 & 4926 & 11 & 2,000000e-02 \\ \hline
scrs8 & 0 & 4477 & 21 & 2,000000e-02 & 0 & 8304 & 21 & 3,000000e-02 \\ \hline
scsd1 & 0 & 1392 & 8 & 1,000000e-02 & 0 & 1472 & 8 & 0,000000e+00 \\ \hline
scsd6 & 0 & 2545 & 11 & 2,000000e-02 & 0 & 2740 & 11 & 2,000000e-02 \\ \hline
scsd8 & 0 & 5879 & 10 & 3,000000e-02 & 0 & 5942 & 10 & 3,000000e-02 \\ \hline
sctap1 & 0 & 2435 & 14 & 1,000000e-02 & 0 & 5367 & 14 & 2,000000e-02 \\ \hline
sctap2 & 0 & 11736 & 12 & 3,000000e-02 & 0 & 55394 & 12 & 1,100000e-01 \\ \hline
sctap3 & 0 & 16100 & 14 & 6,000000e-02 & 0 & 102574 & 14 & 2,600000e-01 \\ \hline
seba & 0 & 53711 & 13 & 1,700000e-01 & 0 & 55513 & 13 & 4,000000e-01 \\ \hline
share1b & 0 & 1417 & 18 & 1,000000e-02 & 0 & 1762 & 18 & 2,000000e-02 \\ \hline
share2b & 0 & 1041 & 17 & 1,000000e-02 & 0 & 1542 & 17 & 1,000000e-02 \\ \hline
shell & 0 & 3983 & 20 & 2,000000e-02 & 0 & 9539 & 20 & 4,000000e-02 \\ \hline
ship04l & 0 & 2641 & 12 & 2,000000e-02 & 0 & 7944 & 12 & 2,000000e-02 \\ \hline
ship04s & 0 & 1778 & 12 & 1,000000e-02 & 0 & 4091 & 12 & 2,000000e-02 \\ \hline
ship08l & 0 & 4442 & 14 & 3,000000e-02 & 0 & 15947 & 14 & 4,000000e-02 \\ \hline
ship08s & 0 & 2295 & 11 & 2,000000e-02 & 0 & 4290 & 11 & 2,000000e-02 \\ \hline
ship12l & 0 & 5506 & 15 & 4,000000e-02 & 0 & 16980 & 15 & 5,000000e-02 \\ \hline
ship12s & 0 & 2507 & 12 & 2,000000e-02 & 0 & 3801 & 12 & 2,000000e-02 \\ \hline
sierra & 0 & 12862 & 18 & 6,000000e-02 & 0 & 43552 & 18 & 1,400000e-01 \\ \hline
stair & 0 & 14682 & 13 & 3,000000e-02 & 0 & 15047 & 13 & 5,000000e-02 \\ \hline
standata & 0 & 2395 & 13 & 1,000000e-02 & 0 & 3953 & 13 & 2,000000e-02 \\ \hline
standgub & 0 & 2395 & 13 & 1,000000e-02 & 0 & 3953 & 13 & 1,000000e-02 \\ \hline
standmps & 0 & 3957 & 23 & 2,000000e-02 & 0 & 9136 & 23 & 3,000000e-02 \\ \hline
stocfor1 & 0 & 805 & 11 & 0,000000e+00 & 0 & 1388 & 11 & 1,000000e-02 \\ \hline
stocfor2 & 0 & 22841 & 19 & 7,000000e-02 & 0 & 32610 & 19 & 2,400000e-01 \\ \hline
tuff & 0 & 7051 & 18 & 2,000000e-02 & 0 & 9820 & 18 & 4,000000e-02 \\ \hline
vtp,base & 0 & 505 & 10 & 0,000000e+00 & 0 & 504 & 10 & 0,000000e+00 \\ \hline
wood1p & 0 & 11645 & 22 & 3,400000e-01 & 0 & 14537 & 22 & 3,600000e-01 \\ \hline
woodw & 0 & 30027 & 30 & 1,800000e-01 & 0 & 126724 & 30 & 5,800000e-01 \\ \hline
\end{tabular}

\end{table}
\begin{table}
  \centering
  \caption{Resultados experimentais da biblioteca Netlib-QAP.}
  \label{tab:resulnetqap}
  \begin{tabular}{|l|r|r|r|r|r|r|r|r|}
\hline
\multicolumn{1}{|c|}{Problema} & \multicolumn{4}{|c|}{MMD} &         \multicolumn{4}{|c|}{RCM} \\ \hline
Nome & R & NNZ & IT & T & R & NNZ & IT & T \\ \hline
qap12 & 15 & 2138580 & 41 & 3,427000e+01 & 15 & 3439794 & 42 & 8,677000e+01 \\ \hline
qap15 & 15 & 8197968 & 46 & 3,353400e+02 & 15 & 13674441 & 54 & 8,185200e+02 \\ \hline
qap8 & 0 & 193944 & 8 & 2,400000e-01 & 0 & 273509 & 32 & 1,720000e+00 \\ \hline
\end{tabular}

\end{table}
\begin{table}
  \centering
  \caption{Resultados experimentais da biblioteca PDS.}
  \label{tab:resulpds}
  \begin{tabular}{|l|r|r|r|r|r|r|r|r|}
\hline
\multicolumn{1}{|c|}{Problema} & \multicolumn{4}{|c|}{MMD} &         \multicolumn{4}{|c|}{RCM} \\ \hline
\multicolumn{1}{|c|}{Nome} & \multicolumn{1}{|c|}{R} &
        \multicolumn{1}{|c|}{NNZ} & \multicolumn{1}{|c|}{IT} &
        \multicolumn{1}{|c|}{T} & \multicolumn{1}{|c|}{R} &
        \multicolumn{1}{|c|}{NNZ} & \multicolumn{1}{|c|}{IT} &
        \multicolumn{1}{|c|}{T} \\ \hline
pds-20 & 0 & 4,1662E+07 & 39 & 2,8013E+03 & 0 & 4,1662E+07 & 39 & 2,8013E+03 \\ \hline
pds-30 & 0 & 1,4644E+07 & 48 & 5,4419E+02 & 0 & 1,4644E+07 & 48 & 5,4419E+02 \\ \hline
pds-40 & 0 & 2,8195E+07 & 50 & 1,6533E+03 & 0 & 2,8195E+07 & 50 & 1,6533E+03 \\ \hline
pds-50 & 0 & 4,1767E+07 & 54 & 3,0229E+03 & 0 & 4,1767E+07 & 54 & 3,0229E+03 \\ \hline
pds-60 & 0 & 5,8119E+07 & 53 & 5,1108E+03 & 0 & 5,8119E+07 & 53 & 5,1108E+03 \\ \hline
pds-70 & -1 & 0,0000E+00 & 0 & 0,0000E+00 & -1 & 0,0000E+00 & 0 & 0,0000E+00 \\ \hline
pds-80 & -1 & 0,0000E+00 & 0 & 0,0000E+00 & -1 & 0,0000E+00 & 0 & 0,0000E+00 \\ \hline
pds-90 & -1 & 0,0000E+00 & 0 & 0,0000E+00 & -1 & 0,0000E+00 & 0 & 0,0000E+00 \\ \hline
pds-100 & -1 & 0,0000E+00 & 0 & 0,0000E+00 & -1 & 0,0000E+00 & 0 & 0,0000E+00 \\ \hline
\end{tabular}

\end{table}
\begin{table}
  \centering
  \caption{Resultados experimentais da biblioteca Rail.}
  \label{tab:resulrail}
  \begin{tabular}{|l|r|r|r|r|r|r|r|r|}
\hline
\multicolumn{1}{|c|}{Problema} & \multicolumn{4}{|c|}{MMD} &         \multicolumn{4}{|c|}{RCM} \\ \hline
\multicolumn{1}{|c|}{Nome} & \multicolumn{1}{|c|}{R} &
        \multicolumn{1}{|c|}{NNZ} & \multicolumn{1}{|c|}{IT} &
        \multicolumn{1}{|c|}{T} & \multicolumn{1}{|c|}{R} &
        \multicolumn{1}{|c|}{NNZ} & \multicolumn{1}{|c|}{IT} &
        \multicolumn{1}{|c|}{T} \\ \hline
rail2586 & 0 & 1,2368E+06 & 90 & 2,3844E+02 & 0 & 1,4425E+06 & 90 & 2,6242E+02 \\ \hline
rail4284 & 0 & 5,8048E+06 & 71 & 6,3984E+02 & 0 & 7,0453E+06 & 71 & 8,4889E+02 \\ \hline
rail507 & 0 & 6,8847E+04 & 36 & 3,3400E+00 & 0 & 7,7942E+04 & 30 & 3,6300E+00 \\ \hline
rail516 & 0 & 5,7220E+04 & 29 & 2,1600E+00 & 0 & 7,3268E+04 & 29 & 2,3700E+00 \\ \hline
rail582 & 0 & 8,6529E+04 & 35 & 4,0400E+00 & 0 & 1,1498E+05 & 35 & 4,3800E+00 \\ \hline
\end{tabular}

\end{table}


% Resultados
%     Descrição dos resultados: deve ser clara e objetiva, resumindo os achados
%     principais que serão detalhados em tabelas e figuras.
%     Ilustrações dos resultados: tabelas e figuras são muito importantes; seu
%     número deve ser o menor possível, e elas devem ser construídas com cuidado
%     para incluir todas as informações necessárias com clareza.
%     Tabelas: devem ser numeradas sequencialmente (Tabela 1, Tabela 2, etc).
%     Seu título deve ser informativo, colocado acima e justificado à esquerda.
%     Notas de rodapé (a, b, c...) podem ser colocados diretamente abaixo da
%     mesma.
%     Figuras (fotos, esquemas, gráficos): devem ser numeradas sequencialmente
%     (Figura 1, Figura 2, etc). Seu título deve ser informativo, colocado
%     abaixo e justificado à esquerda, descrevendo o que é mostrado.

% Discussão / Conclusões
%     Descrição dos dados à luz da literatura
%     Descrição de possíveis fontes de erro e seu efeito sobre os dados
%     Se seus experimentos falharam, quais as sugestões para corrigir o
%     problema?
\section{Conclusões}
Concluímos, pelo menos para os problemas testados, que a heurística de mínimo
grau múltiplo é superior à heurística de Cuthill e McKee por gerar menos
elementos não nulos na decomposição de Cholesky (levando em conta o tamanho do
envelope).


% Matéria encaminhada para publicação
%     Quando houver, referir resumos ou artigos científicos publicados ou
%     encaminhados para publicação
% \input{publicacoes@relatorio_parcial.tex}

% Perspectivas de continuidade ou desdobramento do trabalho
%     O projeto foi concluído ou será continuado?

% Outras atividades de interesse universitário
%     Descrever participações em congressos, cursos extra-curriculares, etc
\section{Apresentações em Congressos}
O aluno também teve o trabalho de iniciação pré-selecionado como um dos cinco
melhores para o Prêmio de Iniciação Científica do XLV Simpósio Brasileiro de
Pesquisa Operacional que será apresentado em setembro de 2013.



% Apoio
%     Citar as agências que financiaram o projeto
% Copyright (C) 2012 Raniere Silva
% 
% This file is part of 'CNPq 126874/2012-3'.
% 
% 'CNPq 126874/2012-3' is licensed under the Creative Commons
% Attribution 3.0 Unported License. To view a copy of this license,
% visit http://creativecommons.org/licenses/by/3.0/.
% 
% 'CNPq 126874/2012-3' is distributed in the hope that it will be
% useful, but WITHOUT ANY WARRANTY; without even the implied warranty of
% MERCHANTABILITY or FITNESS FOR A PARTICULAR PURPOSE.

\section{Informa\c{c}\~{o}es}
Este trabalho foi financiado pelo Conselho Nacional de Desenvolvimento Cient\'{i}fico e Tecnol\'{o}gico pelo Processo 126874/2012-3.

Este trabalho \'{e} licenciado sob a Licen\c{c}a Creative Commons Atribui\c{c}\~{a}o 3.0 N\~{a}o Adaptada License. Para ver uma c\'{o}pia desta licen\c{c}a, visite \url{http://creativecommons.org/licenses/by/3.0/}.
\begin{center}
    \includegraphics{figures/cc-by.png}
\end{center}

% % Copyright (C) 2012 Raniere Silva
% 
% This file is part of 'CNPq 126874/2012-3'.
% 
% 'CNPq 126874/2012-3' is licensed under the Creative Commons
% Attribution 3.0 Unported License. To view a copy of this license,
% visit http://creativecommons.org/licenses/by/3.0/.
% 
% 'CNPq 126874/2012-3' is distributed in the hope that it will be
% useful, but WITHOUT ANY WARRANTY; without even the implied warranty of
% MERCHANTABILITY or FITNESS FOR A PARTICULAR PURPOSE.

\section{FAQ}
\begin{enumerate}
    \item \textbf{Justifique o problema de reduzir a largura de banda é
        convertido em encontrar uma renumeração dos vértices do grafo tal que
        a diferença entre os índices seja mínima?}
        \cite{Fernanda:2005:ReordenacaoCCCG}

    \item \textbf{Por que um nós de alta excentricidade produzem bons
        resultados?} \cite{Fernanda:2005:ReordenacaoCCCG}

    \item \textbf{Por que um limitante inferior para a medida a largura de
        banda de $P A P^t$ para qualquer permutação $P$ é dado pelo menor
        inteiro maior ou igual a $D/2$, onde $D$ é o grau máximo de
        qualquer nó do grafo $G(A)$?} \cite{Cuthill:1969:ReducingBandwidth}

        Tomando o nó de grau máximo a maneira de obter a menor largura de
        banda para ele é posicionando-o na diagonal de modo que metade dos nós
        adjacentes esteja de um lado da banda e a outra metade do outro lado.

    \item \textbf{Por que iniciar com um nó de grau mínimo é uma boa escolha?}
        \cite{Cuthill:1969:ReducingBandwidth}

    \item \textbf{Por que numerar os nós em ordem crescente de grau?}
        \cite{Cuthill:1969:ReducingBandwidth}


\end{enumerate}

% % Copyright (C) 2012 Raniere Silva
% 
% This file is part of 'CNPq 126874/2012-3'.
% 
% 'CNPq 126874/2012-3' is licensed under the Creative Commons
% Attribution 3.0 Unported License. To view a copy of this license,
% visit http://creativecommons.org/licenses/by/3.0/.
% 
% 'CNPq 126874/2012-3' is distributed in the hope that it will be
% useful, but WITHOUT ANY WARRANTY; without even the implied warranty of
% MERCHANTABILITY or FITNESS FOR A PARTICULAR PURPOSE.

\section{Implementações}
A seguir enumeramos algumas das implementações disponíveis.

\subsection{C}
\subsubsection{Independente}
David Fritzsche, implementou o RCM para sua tese de mestrado (``Graph
Theoretical Methods for Preconditioners'', defendida em 2004 no
Departamento de Matemática da Bergische Universität Wuppertal. 

Disponível em \url{http://math.temple.edu/~daffi/software/rcm/}.

\subsection{C++}
\subsubsection{The Boost Graph Library (BGL)}
É uma biblioteca para trabalhar com grafos.

Encontra-se disponível em
\url{http://www.boost.org/doc/libs/1_51_0/libs/graph/doc/index.html}

\subsection{Python}
\subsubsection{PyOrder}
É uma biblioteca para reordenamento de matrizes esparsas.

Encontra-se disponível em \url{https://github.com/dpo/pyorder}.

\subsubsection{NetworkX}
É uma biblioteca para trabalhar com grafos.

Encontra-se disponível em \url{http://networkx.lanl.gov/index.html} e a
implementação do RCM em
\url{http://networkx.lanl.gov/examples/algorithms/rcm.html?highlight=cuthill}


% Bibliografia
%     Diversos formatos: definir qual o mais apropriado
%     IMPORTANTE
%     - Não liste se não citar.
%     - Não cite se não listar.
\bibliographystyle{alpha}
\bibliography{../references}
\end{document}
